\documentclass[11pt,a4paper]{article}
\usepackage[T1]{fontenc}
\usepackage{lmodern,amsmath,amssymb,amsthm,booktabs,longtable,enumitem}
\usepackage[margin=27mm]{geometry}
\usepackage{microtype}
\usepackage[hidelinks]{hyperref}
\hypersetup{pdftitle={Mean-adapted interface bounds and joint source estimates for a theta prime-Gamma energy},pdfauthor={Jongmin Choi},pdfsubject={Full-domain near-half bounds with explicitly separated computational hypotheses},pdfkeywords={theta kernel, prime powers, Gamma form, Carleman operator, source Gram}}
\setlength{\parindent}{1.2em}\setlength{\parskip}{0.22em}
\setlength{\emergencystretch}{2em}\allowdisplaybreaks[2]
\numberwithin{equation}{section}
\newtheorem{theorem}{Theorem}[section]
\newtheorem{proposition}[theorem]{Proposition}
\newtheorem{lemma}[theorem]{Lemma}
\newtheorem{corollary}[theorem]{Corollary}
\newtheorem{hypothesis}[theorem]{Hypothesis}
\theoremstyle{remark}\newtheorem{remark}[theorem]{Remark}
\newcommand{\R}{\mathbb R}\newcommand{\C}{\mathbb C}
\newcommand{\E}{\mathcal E}\newcommand{\D}{\mathcal D}
\newcommand{\He}{L^2_{\mathrm{even}}(\mu)}\newcommand{\Le}{L^2_{\mathrm{even}}(dx)}
\newcommand{\tg}{t_\Gamma}\newcommand{\cg}{c_\Gamma}
\newcommand{\norm}[1]{\lVert#1\rVert}\newcommand{\abs}[1]{\lvert#1\rvert}
\newcommand{\ip}[2]{\langle#1,#2\rangle}\newcommand{\dd}{\mathop{}\!\mathrm d}
\newcommand{\Var}{\operatorname{Var}}\newcommand{\supp}{\operatorname{supp}}
\newcommand{\dom}{\operatorname{Dom}}\newcommand{\one}{\mathbf1}
\newcommand{\PP}{P}\newcommand{\J}{J}\newcommand{\V}{V}
\newcommand{\wo}{\mathcal W}\newcommand{\HS}{\mathrm{HS}}
\title{Mean-adapted interface bounds and joint source estimates for a theta prime--Gamma energy}
\author{Jongmin Choi\\[3pt]\small Independent Researcher, Seoul, Korea\\\small ORCID iD: \href{https://orcid.org/0009-0008-7448-514X}{0009-0008-7448-514X}}
\date{November 1, 2026\\[3pt]\small Manuscript v1.0}
\begin{document}\maketitle
\begin{abstract}
We prove two near-half lower bounds for a complete even theta prime--Gamma energy, with different and explicit computational premises. An accepted whole-local-domain positivity theorem at radius $4/5$, used only as an input theorem, implies the global coefficient $493/1000$. The proof removes an exact rank-one mean component from the Gamma interface and charges its signed contribution to the exterior diagonal. Four prime channels, all six overlaps, both infinite prime tails, and six Gamma strips are retained. A second result, conditional on additional full-residual and centered-response enclosures, gives the coefficient $2861149/5760000$. Its analytic improvement comes from a joint input-Gram estimate for a singular Gamma operator, a regular rank-one majorant, and the complete prime source. Touching kernels are controlled by the Carleman norm rather than a Hilbert--Schmidt estimate. We include proofs of the theta normalization, physical bridge, domain identification, lifted coercivity, and signed exterior cap. The response enclosures are not consequences of local positivity and are not recomputed here. The regularized exterior sign, the global half bound, and the Riemann hypothesis remain unresolved.
\end{abstract}

\section{The two conclusions and their hypotheses}
The two estimates in this paper control the interaction between a fixed local interval and its exterior in different ways. In the first, the original weighted mean repairs a rank-one part of the interface. In the second, mixed products in the input Gram operator reduce the complete physical source norm before a local inverse is applied. The resulting constants have different hypotheses.

Put $B=4/5$, $I=(-B,B)$, $O=\R\setminus[-B,B]$, and $c(x)=\cosh(x/2)$. Define, for $x\ge0$,
\begin{equation}\label{eq:mainkappa}
 \kappa(x)=e^{x/2}\sum_{j\ge1}(4t_j^2-6t_j)e^{-t_j},
 \qquad t_j=\pi j^2e^{2x},
\end{equation}
and extend evenly. Set
\begin{equation}\label{eq:weights}
 W=\frac{\kappa}{2c},\qquad d\mu=2c\kappa\,dx,\qquad
 \rho(s)=\frac{e^{-s/2}}{1-e^{-2s}},\quad
 a_n=\frac{\Lambda(n)}{\sqrt n},\quad \ell_n=\log n.
\end{equation}
Here $\Lambda(p^k)=\log p$ and $\Lambda(n)=0$ otherwise. The measure has mass one. For even profiles let
\begin{align}
 \E[f]={}&\int_0^\infty\rho(s)\int_\R\kappa(x)\kappa(x+s)
 |f(x+s)-f(x)|^2\,dx\,ds\notag\\
 &+\sum_{n\ge2}a_n\int_\R\kappa(x)\kappa(x+\ell_n)
 |f(x+\ell_n)-f(x)|^2\,dx.\label{eq:mainenergy}
\end{align}
Every positive shift is counted once, and every prime power occurs. The domain is
\[
 \D=\{f\in\He:\E[f]<\infty\},\qquad
 \Var_\mu(f)=\norm f_\mu^2-|\mu(f)|^2.
\]
Appendices~\ref{sec:energy}--\ref{sec:domain} prove the equality with the smooth minimum closure and the exact physical bridge. These are analytic conclusions, not additional numerical hypotheses.

For physical numerators put
\begin{align}
 \tg[u]&=\int_0^\infty\rho(s)\norm{u(\cdot+s)-u}_2^2\,ds,
 \label{eq:physicalgamma}\\
 p[u]&=2\sum_{n\ge2}a_n\operatorname{Re}\int_\R
 u(x+\ell_n)\overline{u(x)}\,dx,\label{eq:physicalprime}\\
 q_0[u]&=\tg[u]-\cg\norm u_2^2-p[u],\qquad
 q[u]=q_0[u]+2|m(u)|^2,\quad m(u)=\int_\R cu,\label{eq:qzero}\\
 \cg&=\log\pi-\psi(1/4)=\log\pi+\gamma_E+\pi/2+3\log2.\notag
\end{align}
The local minimum domain $\D_I$ is the closure of even $C_c^\infty(I)$, zero extended, in $(\tg+\norm{\cdot}_2^2)^{1/2}$. Its mean-inclusive closed form $q_I=q|_{\D_I}$ defines a self-adjoint operator $L$.

\begin{hypothesis}[Accepted local input theorem]\label{hyp:local}
On the whole domain $\D_I$, $q_I[u]\ge0$.
\end{hypothesis}
The accepted computer-assisted input states the stronger inequality $q_I[u]\ge10^{-22}\norm u_2^2$. Only its sign is used in the first theorem. Its physical $L^2$ coefficient is not transferred unchanged to a weighted global gap. Appendix~\ref{app:provenance} identifies the accepted evidence; it does not claim a fresh replay.

\begin{theorem}[Local-sign branch]\label{thm:localglobal}
Under Hypothesis~\ref{hyp:local}, every even $f\in\D$ satisfies
\begin{equation}\label{eq:493}
 \E[f]\ge\frac{493}{1000}\Var_\mu(f).
\end{equation}
The same comparison permits the endpoint
\begin{equation}\label{eq:lambdaendpoint}
 \lambda_*=\frac{19985-\sqrt{69921}}{40000}>\frac{493}{1000},
 \qquad \E[f]\ge\lambda_*\Var_\mu(f).
\end{equation}
No local inverse-response enclosure is needed.
\end{theorem}

\begin{theorem}[Additional-response branch]\label{thm:conditionalglobal}
Under the complete-operator and centered-response contracts in Hypothesis~\ref{hyp:response}, every even $f\in\D$ satisfies
\begin{equation}\label{eq:newglobal}
 \E[f]\ge\frac{2861149}{5760000}\Var_\mu(f).
\end{equation}
These contracts imply local positivity but are not implied by Hypothesis~\ref{hyp:local}. Their numerical enclosure assertions are inputs, not computations established here.
\end{theorem}
The latter coefficient exceeds $71501077/144000000$, obtained from a larger source cost under the same contracts, by exactly $24/125000$. This comparison concerns conditional implications, not the endpoint $1/2$. Arguments are given for real-even functions; decomposition into real and imaginary parts gives the complex statements. Inner products are linear in the first argument. This is a research manuscript, not an accepted or peer-reviewed publication. No historical-priority or formal-verification claim is made.

\section{Analytic estimates and the two source coordinates}\label{sec:estimates}
Write $\alpha=\sqrt W$, and set
\[
 A(x)=\frac{(4t^2-6t)e^{-t}}{1+e^{-x}}+128t^2e^{-4t},
 \quad t=\pi e^{2x},\qquad A=A(B).
\]
Appendix~\ref{app:scalar} proves
\begin{gather}
 A<1/9216,\quad W(B+h)\le Ae^{-25h},\quad
 \int_B^\infty W<1/230400,\label{eq:envelope}\\
 \sup W<1/2,\quad \int_0^\infty W\le1/8,\quad
 p_*:=\mu(O)<3/49000,\quad \cg<383/70<11/2,\label{eq:basiccaps}\\
 W(1/10)<21/50,\quad W(1/4)<1/4,\quad W(9/20)<3/40,\label{eq:pointW}\\
 \kappa(2/5)<1/4,\quad \kappa(1/2)<7/100,\quad
 \kappa(3/5)<1/50,\quad \kappa(7/10)<3/1000,\quad
 \kappa(3/4)<1/1000.\label{eq:pointkappa}
\end{gather}
Both $\kappa$ and $W$ decrease on $[1/10,\infty)$. Also,
\begin{equation}\label{eq:rhoelementary}
 \rho(s)\le\frac1{2s}+\frac14,\qquad
 \frac56<\rho(B)<\frac9{10},\qquad \pi^2<10.
\end{equation}
These are full-series analytic inequalities.

There are two source norms. The weighted profile interface $\mathcal J=\mathcal J_p+\mathcal J_\Gamma$ is the inside-to-outside operator in
\begin{equation}\label{eq:energycross}
 \E[g+h]=\E[g]+\E[h]-2\ip{\mathcal Jg}{h}_\mu,
 \qquad \supp g\subset\overline I,\quad\supp h\subset\overline O.
\end{equation}
Under the unitary map
\begin{equation}\label{eq:foldU}
 (Uf)(x)=\sqrt{4c(x)\kappa(x)}f(x),\qquad x>0,
\end{equation}
its Gamma kernel is $\sqrt{W(x)W(y)}[\rho(x-y)+\rho(x+y)]$, $x>B$, $0<y<B$.

The physical exterior source in the second branch is
\begin{equation}\label{eq:physicalsource}
 (Ju)(y)=\int_O\rho(|x-y|)u(x)\,dx+
 \sum_{n\ge2}a_n[u(y+\ell_n)+u(y-\ell_n)],\qquad y\in I,
\end{equation}
where exterior inputs are extended by zero. Its input norm is
$\wo(u)=\int_O|u|^2/W\,dx$, and its output norm is physical $L^2(I)$.
With $v(x)=\sqrt2u(x)/\alpha(x)$ on $x>B$, the input norm becomes $\norm v_2^2=\wo(u)$ and the output fold is multiplication by $\sqrt2$. The extra factor $\sqrt{W(y)}$ belongs only to the profile interface.

\begin{lemma}[Carleman bound]\label{lem:carleman}
The kernel $(s+t)^{-1}$ on the positive half-line has norm at most $\pi$, as does every rectangular restriction.
\end{lemma}
\begin{proof}
For the Schur weight $w(s)=s^{-1/2}$,
\[
 \int_0^\infty\frac{t^{-1/2}}{s+t}\,dt
 =2s^{-1/2}\int_0^\infty\frac{dr}{1+r^2}=\pi s^{-1/2}.
\]
The transposed integral is identical. Weighted Cauchy--Schwarz proves the assertion on compact tests away from zero, and density completes it.
\end{proof}

\section{Four weighted prime channels and all their overlaps}\label{sec:weightedprime}
For the folded adjoint of the profile interface the branches are
\begin{align*}
 (P_nv)(y)&=a_n\sqrt{W(y)W(\ell_n+y)}v(\ell_n+y)\one_{\ell_n+y>B},\\
 (M_nv)(y)&=a_n\sqrt{W(y)W(\ell_n-y)}v(\ell_n-y)\one_{\ell_n-y>B}.
\end{align*}
If $y-\ell_n>0$, it is still inside; a negative exterior input gives exactly the reflected-minus branch. Retain $H_p=P_2+M_3+M_4+M_5$. Its domains, up to null endpoints, are
\begin{center}\small
\begin{tabular}{@{}lll@{}}\toprule
Channel & Inside output $y$ & Exterior input $x$\\\midrule
$P_2$ & $(B-\log2,B)$ & $(B,B+\log2)$\\
$M_3$ & $(0,\log3-B)$ & $(B,\log3)$\\
$M_4$ & $(0,\log4-B)$ & $(B,\log4)$\\
$M_5$ & $(0,B)$ & $(\log5-B,\log5)$\\\bottomrule
\end{tabular}
\end{center}
Here $M_2=0$ and $a_6=0$. Appendix~\ref{app:scalar} gives
\begin{equation}\label{eq:coefficientcaps}
 a_2<99/200,\quad a_3<16/25,\quad a_4<7/20,\quad a_5<29/40,
 \qquad a_n\le1.
\end{equation}
In particular $a_4=(\log2)/2$, not $(\log4)/2$.

The four channel-square weights, in units of $A$, are bounded by
$21/50,1/4,3/40,1/100$. For $P_2$, $y>B-\log2>1/10$.
For $M_3$ split at $y=1/4$: below it $x-B>6/125$ gives
$W(x)W(y)<Ae^{-6/5}/2<A/4$.
For $M_4$ split at $9/20$: below it $x-B>17/125$ gives $Ae^{-17/5}/2<A/40$.
For $M_5$ split at $3/5$: below it $x-B>209/1000$ gives $A/280<A/100$.
Above each split use \eqref{eq:pointW}--\eqref{eq:pointkappa}.
The common input diagonal is at most
\begin{align}
 A\bigl[(99/200)^2(21/50)&+(16/25)^2/4
 +(7/20)^2(3/40)+(29/40)^2/100\bigr]\notag\\
 &=A\frac{879017}{4000000}<\frac{11A}{50}.\label{eq:weighteddiag}
\end{align}
No output orthogonality is assumed.

A cross operator has the form $Cv(x)=\one_D(x)w(x)v(\phi(x))$,
where $\phi$ is injective and $|\phi'|=1$, so $\norm C\le\sup_D|w|$.
The pair $C+C^*$ costs twice this bound, even when $D$ is not reflection invariant.
Every pair is listed below; the last column bounds $\sqrt{W(x)W(x')}/A$.
\begin{center}\small
\begin{tabular}{@{}llll@{}}\toprule
Pair & First input domain & Second input $x'$ & Cap\\\midrule
$P_2,M_3$ & $(B,\log6-B)$ & $\log6-x$ & $1/10$\\
$P_2,M_4$ & $(B,\log8-B)$ & $\log8-x$ & $1/300$\\
$P_2,M_5$ & $(B,B+\log2)$ & $\log10-x$ & $1/6000$\\
$M_3,M_4$ & $(B,\log3)$ & $x+\log(4/3)$ & $1/32$\\
$M_3,M_5$ & $(B,\log3)$ & $x+\log(5/3)$ & $1/500$\\
$M_4,M_5$ & $(B,\log4)$ & $x+\log(5/4)$ & $1/15$\\\bottomrule
\end{tabular}
\end{center}
These rational comparisons follow from the envelope and are proved in Appendix~\ref{app:scalar}. In the first three rows use $W(y)<21/50$, and in the next two use $W(y)<1/2$.
In the last row the sharper full product is below $A/200$:
for $y\ge9/20$ use $(3/40)(1/15)$; otherwise the sum of exterior distances is
$\log20-2B-2y>495/1000$, its exponent exceeds $6$, and the product is below $A/800$.
The six paired costs in units of $A$ are at most
\[
 27/1000,\quad1/2000,\quad3/50000,\quad7/1000,\quad1/1000,\quad13/5000.
\]
Their sum is $1908/50000<1/25$. Thus $\norm{H_p}^2<(13/50)A<(51/100)^2A$.

The exact complement is $\sum_{n\ge3}P_n+\sum_{n\ge7}M_n$.
From $e^B<9/4$, $e^{2B}<5$ and $25/2>12$,
\begin{gather}
 \norm{P_n}\le\sqrt{A/2}(9/(4n))^{12},\qquad
 \norm{M_n}\le\sqrt{A/2}(5/n)^{12},\notag\\
 \sum_{n\ge3}(9/(4n))^{12}\le\frac{14}{11}(3/4)^{12}<1/24,\qquad
 \sum_{n\ge7}(5/n)^{12}\le\frac{18}{11}(5/7)^{12}<3/100.\label{eq:bothtails}
\end{gather}
The last comparisons are $178564176<184549376$ and
$439453125000<456762477633$. Both complete tails converge in operator norm.
Their triangle inequalities pay all head--tail and tail--tail products.
Since $\sqrt2>7/5$,
\begin{equation}\label{eq:weightedprimenorm}
 \norm{\mathcal J_p}<\sqrt A\left(\frac{51}{100}+\frac{43}{840}\right)
 <\frac9{16}\sqrt A<\frac3{512}<\frac{59}{10000}.
\end{equation}
The scalar slack before multiplication by $\sqrt A$ is $11/8400$.

\section{The exact mean repair and six Gamma strips}\label{sec:meanrepair}
For $x>B$ and $0<y<B$ define
\begin{equation}\label{eq:remainderkernel}
 B_\Gamma(x,y)=\rho(x-y)+\rho(x+y)-2c(y)\rho(x).
\end{equation}
The full positive series for $\rho$ gives
\[
 B_\Gamma(x,y)=2\sum_{j\ge1}e^{-(2j+1/2)x}
 [\cosh((2j+1/2)y)-\cosh(y/2)]\ge0.
\]
It decreases in $x$ and increases in $y$. Only the $j=0$ term cancels.
Let $T_\Gamma$ have kernel $\sqrt{W(x)W(y)}B_\Gamma(x,y)$ and set
\[
 e(y)=U1(y)=2\sqrt2c(y)\sqrt{W(y)},\qquad
 z(x)=\frac{\sqrt{W(x)}\rho(x)}{\sqrt2}.
\]
For an inside profile $g$ the exact identity is
\begin{equation}\label{eq:rankrepair}
 U\mathcal J_\Gamma g=T_\Gamma Ug+z\mu(g),\qquad
 \mu(g)=\ip{Ug}{e}.
\end{equation}
Thus the rank-one term has not been discarded.

For $y\le b<B$, monotonicity and \eqref{eq:rhoelementary} give
\[
 0\le B_\Gamma(x,y)\le B_\Gamma(B,b)
 <\frac{B}{B^2-b^2}-\frac76=:F(b).
\]
The subtraction is made at $x=B$, not by imposing a false lower bound for $\rho(x)$ in the exterior.
Divide the input at $2/5,1/2,3/5,7/10,3/4$.
The first five strips give valid Hilbert--Schmidt estimates:
\begin{center}\small
\begin{tabular}{@{}llll@{}}\toprule
Strip & $\int W(y)\,dy$ cap & $F(b)$ & Squared norm cap\\\midrule
$(0,2/5)$ & $1/8$ & $1/2$ & $1/7372800$\\
$(2/5,1/2)$ & $1/80$ & $23/26$ & $529/12460032000$\\
$(1/2,3/5)$ & $7/2000$ & $71/42$ & $5041/116121600000$\\
$(3/5,7/10)$ & $1/1000$ & $25/6$ & $625/8294400000$\\
$(7/10,3/4)$ & $3/40000$ & $1703/186$ & $2900209/106278912000000$\\\bottomrule
\end{tabular}
\end{center}
Each row uses the exterior integral $1/230400$.
The squared norms are respectively below $(140,43,44,76,28)/10^9$.

On the touching strip $(3/4,4/5)$, $W(y)<1/2000$ and
\[
 B_\Gamma(x,y)\le\rho(x-y)+\rho(x+y)
 \le\frac1{2(x-y)}+\frac{51}{62}.
\]
The singular term is half a Carleman kernel after $s=x-B$, $t=B-y$.
Its norm is below
$\frac\pi2(9216\cdot2000)^{-1/2}<37/100000$.
Its squared cap is $1/7372800$, and $1369\cdot7372800>10^{10}$.
The bounded term has norm at most
\[
 \frac{51}{62}(230400\cdot40000)^{-1/2}
 =\frac{17}{1984000}<\frac9{1000000}.
\]
Thus the last strip norm is below $379/10^6<19/50000$, whose square is below $145/10^9$.
No Hilbert--Schmidt claim is made for its touching singularity.
The input projections are orthogonal, giving
$T_\Gamma T_\Gamma^*=\sum_iT_iT_i^*$, without output orthogonality.
Consequently
\begin{equation}\label{eq:gammaremainnorm}
 \norm{T_\Gamma}^2<\frac{140+43+44+76+28+145}{10^9}
 =\frac{476}{10^9}<\left(\frac7{10000}\right)^2.
\end{equation}
Both physical boundaries are already included in the even fold.

\section{The global comparison from local sign}\label{sec:firstproof}
The complete exterior prime operator has norm below $(4/5)A$.
There are two same-side translations and one opposite-side reflection; they cost
\[
 A\left[2\sum_{n\ge2}\Lambda(n)n^{-13}
 +e^{20}\sum_{n\ge5}\Lambda(n)n^{-13}\right].
\]
Using $\Lambda(n)\le n/2$, isolating $n=5$ with $a_5<3/4$, and retaining every $n\ge7$ term gives the bracket
\[
 \frac{13}{11\cdot2^{12}}+\frac34+
 \frac98\,5^{12}\frac{18}{11}7^{-12}<\frac45,
\]
with slack $4887143214909/283469561876480$; $a_6=0$ is the only intervening omission.
The factors use $e^{20}<5^{25/2}$ and $e^{20}<\frac94\,5^{12}$, both consequences of $e^{8/5}<5$ and $e^{4/5}<9/4$.
This is an exterior operator estimate, not a reverse-interface row estimate.

For exterior $h$, the physical bridge gives
\[
 \E[h]=\tfrac12\norm h_\mu^2+\tg[\kappa h]-\cg\norm{\kappa h}_2^2
 -\ip{K_{p,OO}h}{h}_\mu.
\]
Since $\norm{\kappa h}_2^2\le A\norm h_\mu^2$,
\begin{equation}\label{eq:exteriorenergy}
 \E[h]\ge(\tfrac12-\varepsilon)\norm h_\mu^2,\qquad
 \varepsilon=\frac{11/2+4/5}{9216}=\frac7{10240}.
\end{equation}
Also
\[
 \norm z^2<\tfrac12(9/10)^2/230400,\qquad
 \norm z<\frac3{1600\sqrt2}.
\]
For centered $v=g+h$, $\mu(g)=-\mu(h)$ and
$|\mu(h)|\le\sqrt{p_*}\norm h_\mu$.
The removed cross term is
$-2\mu(g)\ip z{Uh}=2\mu(h)\ip z{Uh}$.
It is a product of distinct exterior functionals and is not assigned a sign.
Its absolute charge is bounded by
\begin{equation}\label{eq:rankcharge}
 2\sqrt{p_*}\norm z<\frac3{800}\sqrt{\frac3{98000}}<\frac1{48000},
\end{equation}
where squaring reduces the last comparison to $62208000000<62720000000$.

Local positivity gives
$\E[g]\ge\frac12\norm g_\mu^2-\frac12|\mu(g)|^2$.
Together with the complete interface bounds this yields
\begin{equation}\label{eq:twobytwo}
 \E[v]\ge\tfrac12\norm g_\mu^2+(\tfrac12-a)\norm h_\mu^2
 -2L_0\norm g_\mu\norm h_\mu,\qquad
 a=\frac3{4000},\quad L_0=\frac{33}{5000}.
\end{equation}
Indeed the cross constant is below $59/10000+7/10000=L_0$, and
\[
 \varepsilon+p_*/2+1/48000
 <7/10240+1/32000+1/48000=113/153600<3/4000.
\]
The final slack is $11/768000$.

At $\delta=7/1000$ the residual comparison matrix is
\[
 M_\delta=\begin{pmatrix}7/1000&-33/5000\\-33/5000&1/160\end{pmatrix},
 \quad \det M_\delta=19/10^8,\quad
 \operatorname{tr}M_\delta=53/4000.
\]
It is positive definite, and
\begin{equation}\label{eq:margin}
 M_\delta\succeq\frac{\det M_\delta}{\operatorname{tr}M_\delta}I
 =\frac{19}{1325000}I.
\end{equation}
More generally it is positive semidefinite for
$\delta\ge(15+\sqrt{69921})/40000$, proving the stated endpoint.
At $\delta=1/200$ its determinant is $-2231/10^8$.
Thus this comparison does not certify $99/200$ or $1/2$; its failure does not imply a negative direction of the actual form.

First apply the argument to $v=f-\mu(f)$ with $f$ even compact smooth.
Constants and sharp cuts belong to the corresponding minimum domains by Appendix~\ref{app:cuts}.
The bounded interface operators and norm-convergent tails justify the cross identity on these cuts.
Finally use the global smooth form core, along which energy and variance converge.
This proves Theorem~\ref{thm:localglobal} on the full maximal domain.

\section{A joint Gram estimate for the physical source}\label{sec:joint}
The output coordinate in this section is physical $L^2(I)$, not weighted profile space.
For $v(x)=\sqrt2u(x)/\alpha(x)$ on $x>B$ the folded Gamma source is
\[
 (T_\Gamma^{\rm phys}v)(y)=\int_B^\infty
 [\rho(x-y)+\rho(x+y)]\alpha(x)v(x)\,dx,\qquad0<y<B.
\]
There is no further factor $\sqrt2$.
Because $x+y>B$, \eqref{eq:rhoelementary} implies the positive majorant
\[
 0\le\rho(x-y)+\rho(x+y)\le\frac1{2(x-y)}+b,\qquad b=\frac{23}{20}.
\]
Define
\[
 (Cv)(y)=\int_B^\infty\frac{\alpha(x)v(x)}{2(x-y)}\,dx,\qquad
 (Rv)(y)=b\int_B^\infty\alpha(x)v(x)\,dx.
\]
Pointwise $|T_\Gamma^{\rm phys}v|\le(C+R)|v|$.
This is domination, not equality with the actual Gamma source.
The touching kernel is half a Carleman kernel, hence
\[
 \norm C\le\frac\pi2\sqrt A.
\]
Set
\[
 I_0=\int_B^\infty\alpha(x)^2\,dx,\qquad
 I_2=\int_B^\infty\alpha(x)^2\log^2\frac{x}{x-B}\,dx.
\]
Then $I_0\le A/25$ and
\begin{equation}\label{eq:regulargram}
 \norm R^2=b^2BI_0\le\frac{529}{12500}A,\qquad
 \norm{C^*R}=\frac b2\sqrt{I_2I_0}.
\end{equation}
The second identity follows from the exact rank-one formula
\[
 (C^*Rv)(x)=\frac b2\alpha(x)\log\frac{x}{x-B}
                  \int_B^\infty\alpha(t)v(t)\,dt.
\]

\subsection{The logarithmic cross integral}
With $t=25(x-B)$,
\[
 I_2/A\le\frac1{25}\int_0^\infty e^{-t}\log^2(1+20/t)\,dt.
\]
Take $a=61/20>\log21$, since $e^3>20$ and $e^{1/20}>21/20$.
Writing $r(t)=\max\{-\log t,0\}$, one has
$\log(1+20/t)\le a+r(t)$.
The inequality $e^{-t}\le1-t+t^2/2$ on $[0,1]$ and
\[
 \int_0^1t^k(-\log t)\,dt=(k+1)^{-2},\qquad
 \int_0^1t^k\log^2t\,dt=2(k+1)^{-3}
\]
give
\[
 \int_0^1e^{-t}(-\log t)\,dt\le29/36,\qquad
 \int_0^1e^{-t}\log^2t\,dt\le193/108.
\]
Consequently
\begin{equation}\label{eq:I2}
 I_2/A\le\frac1{25}\left(\frac{3721}{400}+
       \frac{1769}{360}+\frac{193}{108}\right)
 =\frac{172837}{270000}<\left(\frac{81}{100}\right)^2.
\end{equation}
Thus $2\norm{C^*R}<(1863/10000)A$ and
\begin{equation}\label{eq:gammajoint}
 \norm{J_\Gamma}_{\wo\to L^2(I)}^2/A
 <\frac{121}{49}+\frac{529}{12500}+\frac{1863}{10000}
 <\frac{529}{196}.
\end{equation}
The last margin is $9524/9800000$.
This proves the standalone bound
$\norm{J_\Gamma}_{\wo\to L^2(I)}<(23/14)\sqrt A$.
The calculation concerns $C^*R$, not the square integral of a touching Gamma kernel.

\subsection{The physical prime head and complete tails}
In the same folded spaces the physical channels are
\[
 (P_nv)(y)=a_n\alpha(\ell_n+y)v(\ell_n+y)\one_{\ell_n+y>B},\quad
 (M_nv)(y)=a_n\alpha(\ell_n-y)v(\ell_n-y)\one_{\ell_n-y>B}.
\]
In this section $H_p=P_2+M_3+M_4+M_5$ refers to these physical channels.
They have the same domains and pair maps as in Section~\ref{sec:weightedprime}, but without $\sqrt{W(y)}$.
Their common input diagonal is at most $(422091/400000)A$.
Indeed below $x_5=\log5-B$ the square sum is at most $(6217/8000)A$;
above $x_5$, $\alpha(x)/\sqrt A<e^{-9/80}<9/10$, giving
\[
 D(x)/A\le\frac{81}{100}\bigl[(99/200)^2+(16/25)^2
 +(7/20)^2+(29/40)^2\bigr]=\frac{422091}{400000}.
\]
The six cross products use the last column of the table in Section~\ref{sec:weightedprime}.
Their paired sum is
\[
 \frac{198}{3125}+\frac{231}{200000}+\frac{957}{8000000}
 +\frac7{500}+\frac{116}{62500}+\frac{203}{6000}
 =\frac{109751}{960000}.
\]
Hence
\begin{equation}\label{eq:physicalhead}
 \norm{H_p}^2\le\frac{5613847}{4800000}A<\frac{169}{144}A.
\end{equation}
The exact complement $T_p=\sum_{n\ge3}P_n+\sum_{n\ge7}M_n$ obeys
\begin{equation}\label{eq:physicaltail}
 \norm{T_p}<\left(\frac1{24}+\frac3{100}\right)\sqrt A
\end{equation}
by the same complete integer tails \eqref{eq:bothtails}, now without $1/\sqrt2$.
Thus no prime power or head--tail cross term is omitted.

\subsection{Every singular Gamma and prime-head product}
Let $\sigma=25/2$. Direct substitution gives the \emph{input}-Gram kernels
\begin{align}
 (C^*P_n)(z,t)&=\frac{a_n\alpha(z)\alpha(t)}
                      {2(z-t+\ell_n)},&
 z>B,\quad\max(B,\ell_n)<t<B+\ell_n,\label{eq:inputgramplus}\\
 (C^*M_n)(z,t)&=\frac{a_n\alpha(z)\alpha(t)}
                      {2(z+t-\ell_n)},&
 z>B,\quad\max(B,\ell_n-B)<t<\ell_n.\label{eq:inputgramminus}
\end{align}
They vanish outside the indicated intervals. They are not the output-Gram products $CP_n^*$ or $CM_n^*$.

For $P_2$, put $\ell=\log2$, $s=z-B$, $r=t-B\in(0,\ell)$.
The kernel is bounded by
$(a_2A/2)e^{-\sigma(s+r)}/(s+\ell-r)$.
On $r\le\ell/2$ the denominator is at least $\ell/2$, so an actual Hilbert--Schmidt majorant gives
\[
 \norm{(C^*P_2)_{\rm low}}/A
 \le\frac{a_2}{2\sigma\ell}<\frac{297}{10000}.
\]
On $r>\ell/2$, set $h=\ell-r$. The denominator is $s+h$, so use Carleman, retaining the factor $1/2$:
\[
 \norm{(C^*P_2)_{\rm high}}/A
 \le\frac{a_2\pi}{2}\,2^{-25/4}<\frac{1089}{89600}.
\]
This corner $(z,t)=(B,B+\log2)$ is not Hilbert--Schmidt.

For $M_3,M_4$, let $d_n=2B-\ell_n$. The bounds
$d_3>1/2$, $d_4>1/5$ genuinely separate their kernels.
The exponential square integrals give
\[
 \norm{C^*M_n}/A\le\frac{a_n}{4\sigma d_n},\qquad
 \norm{C^*M_3}/A<16/625,\quad \norm{C^*M_4}/A<7/200.
\]
For $M_5$, however, set $d=\ell_5-2B>0$,
$s=z-B$ and $h=t-(\ell_5-B)$.
The kernel majorant is
\[
 \frac{a_5A}{2}e^{-\sigma d}\frac{e^{-\sigma(s+h)}}{s+h}.
\]
Although its two input locations are separated, the corner is touching.
Carleman, not a denominator gap or Hilbert--Schmidt estimate, gives
$\norm{C^*M_5}/A\le a_5\pi/2<319/280$.
All four products are therefore paid by
\begin{equation}\label{eq:CHp}
 \norm{C^*H_p}/A<
 \frac{297}{10000}+\frac{1089}{89600}+\frac{16}{625}
 +\frac7{200}+\frac{319}{280}
 =\frac{2781497}{2240000}<\frac54.
\end{equation}

\subsection{The regular-majorant product and the full source}
On the common exterior input coordinate,
\begin{align*}
 H_p^*1(t)&=\alpha(t)w(t),\\
 w(t)&=a_2\one_{(B,B+\log2)}(t)+a_3\one_{(B,\log3)}(t)
 +a_4\one_{(B,\log4)}(t)+a_5\one_{(\log5-B,\log5)}(t).
\end{align*}
Since $w<221/100$ and $R=b\,1\otimes\alpha$, the exact rank-one product satisfies
\begin{equation}\label{eq:HpR}
 H_p^*R=b(H_p^*1)\otimes\alpha,\qquad
 2\norm{H_p^*R}\le2b(221/100)I_0\le\frac{5083}{25000}A.
\end{equation}
Expanding the bounded-operator Gram of $C+R+H_p$ now charges all three cross pairs:
\begin{align}
 \norm{C+R+H_p}^2/A
 &<\frac{121}{49}+\frac{169}{144}+\frac52+
   \frac{529}{12500}+\frac{1863}{10000}+\frac{5083}{25000}\notag\\
 &=\frac{121}{49}+\frac{169}{144}+\frac52+\frac{21597}{50000}
 <\frac{324}{49}.\label{eq:jointgram}
\end{align}
For the last comparison use $21597/50000<9/20$ and
$169/144+59/20=2969/720<203/49$, whose cross products are
$145481<146160$.
Positive-kernel domination gives
$|Jv|\le(C+R+H_p+T_p)|v|$.
Combining \eqref{eq:physicaltail} and \eqref{eq:jointgram} proves
\begin{theorem}[Complete analytic source bound]\label{thm:source}
The physical source extends from compact exterior tests to every even finite-$\wo$ input and satisfies
\begin{equation}\label{eq:sourcej}
 j:=\norm J_{\wo\to L^2(I)}
 <\left(\frac{18}{7}+\frac1{24}+\frac3{100}\right)\sqrt A
 =\frac{11101}{4200}\sqrt A<\frac{11101}{403200}<\frac{69}{2500}.
\end{equation}
This theorem needs neither local sign nor response enclosures.
\end{theorem}
The extension follows by density and boundedness; both full prime tails converge in operator norm. The domination statement is not a Loewner order between rectangular operators.

\section{The additional response contracts and inverse cost}\label{sec:response}
Let $e_0=(2B)^{-1/2}$ and
$e_n=B^{-1/2}\cos(n\pi x/B)$, $n\ge1$, on $I$.
Let $E:\R^{64}\to L^2_{\rm even}(I)$ insert $e_0,\ldots,e_{63}$ and set
\[
 H_{64}=L+9EE^*,\qquad U=H_{64}^{-1}E.
\]
Appendix~\ref{sec:lift} proves analytically that
\begin{equation}\label{eq:Hcoercivity}
 H_{64}\succeq61I/28,\qquad H_{64}^{-1}\preceq D_{\rm old},
\end{equation}
where $D_{\rm old}$ is the explicit cosine-diagonal majorant defined there.
These facts are not numerical enclosures.

Choose fixed real trial coefficients $C_{\rm trial}\in\R^{1024\times64}$,
defining columns $\widetilde U$ in the first 1024 local cosine modes.
These columns belong to the local operator domain: their killed Gamma row has logarithmic endpoint singularities in $L^2$, and their remaining rows are bounded.
Set
\[
 R_{64}=E-H_{64}\widetilde U,\qquad
 F=I_{64}/9-E^*H_{64}^{-1}E.
\]
The regular-majorant operator $R$ in Section~\ref{sec:joint} is unrelated to $R_{64}$.

\begin{hypothesis}[Full residual and centered-response contracts]\label{hyp:response}
For the fixed trial matrix there are symmetric matrices $F_*,V_{\rm up}$,
and its positive-diagonal Cholesky factor $M$ with $F_*=MM^*$, such that:
\begin{enumerate}[label=\textup{(\alph*)}]
\item The actual complete-operator inequalities hold:
\begin{equation}\label{eq:contracta}
 0\prec F_*\preceq F,\qquad
 R_{64}^*H_{64}^{-1}R_{64}
 \preceq R_{64}^*D_{\rm old}R_{64}
 \preceq V_{\rm up}\preceq F_*/100.
\end{equation}
All residual rows, including those beyond 1024 and the infinite remainder, are included.
\item Set $V=\widetilde U M^{-*}$ and $Z=C_{\rm trial}M^{-*}$.
For one fixed rational centering vector $a\in\R^{64}$, the exact finite-row, spatial-tail, Gamma-tail, and prime-tail quantities specified in Appendix~\ref{app:rows} have enclosures satisfying
\begin{equation}\label{eq:contractb}
 \sqrt{I_{\rm fin}}+e_X<\frac{81}{2500},\qquad
 e_\Gamma<\frac{39}{10000},\qquad
 b_{16}\norm{Z_{<64}}_{\HS}<10^{-21},\qquad
 \norm{Z_{64:1023}}_{\HS}<\frac13.
\end{equation}
Here $b_{16}=500/20^{25}$ bounds the entire omitted prime source after 16.
The rows and their centering refer to the same $V$, not to independently chosen response vectors.
\end{enumerate}
\end{hypothesis}
The theorem is valid for any fixed data satisfying these precise contracts.
An accepted finite matrix sign or a finite replay alone does not establish their full analytic enclosure interpretation.
Neither \eqref{eq:contracta} nor \eqref{eq:contractb} is inferred from Hypothesis~\ref{hyp:local}.
No new enclosure run is asserted.

To see local coercivity from (a), put $f_*=\lambda_{\min}(F_*)>0$ and
$K_0=3H_{64}^{-1/2}E$. Then
$K_0^*K_0=I-9F\preceq(1-9f_*)I$, so the same upper bound holds for
$K_0K_0^*$. Therefore
\begin{equation}\label{eq:Lcoercive}
 L=H_{64}^{1/2}(I-K_0K_0^*)H_{64}^{1/2}
 \succeq9f_*H_{64}\succeq(549f_*/28)I>0.
\end{equation}
The mean pole remains in $L$ throughout.

For $m_O(u)=\int_Ocu=0$, the fixed centered-row decomposition in Appendix~\ref{app:rows} gives
\begin{equation}\label{eq:responsebound}
 \norm{V^*Ju}\le r\sqrt{\wo(u)},\qquad
 r<\frac{363}{10000}+\frac j3+10^{-21}<\frac{91}{2000}.
\end{equation}
It is essential to insert the sharper fraction from \eqref{eq:sourcej} before rounding:
\[
 \frac{91}{2000}-\frac{363}{10000}-\frac{11101}{1209600}
 =\frac{683}{30240000}>10^{-21}.
\]
The coarser bound $j<69/2500$ alone leaves no displayed room for this last tail.

Woodbury gives the actual inverse
\[
 T=L^{-1}=H_{64}^{-1}+UF^{-1}U^*.
\]
For $s\in L^2_{\rm even}(I)$ set $h=\ip s{H_{64}^{-1}s}$.
Since $U=\widetilde U+H_{64}^{-1}R_{64}$, inverse order and (a) imply
\[
 \ip s{Ts}\le h+
 \left[\norm{V^*s}+\frac1{10}\sqrt h\right]^2.
\]
Indeed
$\norm{M^{-1}R_{64}^*H_{64}^{-1}s}\le\sqrt h/10$
follows by conjugating \eqref{eq:contracta}.
Young's inequality with parameter $1/20$ yields
\[
 \ip s{Ts}\le\frac{121}{100}h+\frac{21}{20}\norm{V^*s}^2.
\]
For $s=Ju$, $h\le(28/61)j^2\wo(u)$.
Using \eqref{eq:responsebound} and then $j<69/2500$,
\begin{align}
 \frac{\ip{Ju}{TJu}}{\wo(u)}
 &<\frac{847}{1525}\left(\frac{69}{2500}\right)^2
 +\frac{21}{20}\left(\frac{91}{2000}\right)^2\notag\\
 &=\frac{4032567}{9531250000}+\frac{173901}{80000000}
 <\frac{17}{40000}+\frac{87}{40000}=\frac{13}{5000}.\label{eq:inversecost}
\end{align}
For the first comparison use $847/1525<5/9$ and
$17/40000-4761/11250000=81/45000000>0$; the second uses $173901<174000$.
Including $u=0$, the universal statement is
$\ip{Ju}{L^{-1}Ju}\le(13/5000)\wo(u)$ on centered exterior inputs.

\section{The signed exterior form and the conditional global bound}\label{sec:conditionalproof}
Let $\D_O^0$ be the minimum closure of even $C_c^\infty(O)$ with
$m_O(u)=0$ in $(\tg+\wo)^{1/2}$.
For $u\in\D_O^0$ define
\[
 S[u]=\tg[u]-\cg\norm u_2^2-p[u]-\ip{Ju}{L^{-1}Ju}.
\]
This is the actual signed Schur form of the mean-inclusive local operator.
Appendix~\ref{sec:signed} proves, without a numerical premise,
\begin{equation}\label{eq:signedcapmain}
 \cg\norm u_2^2+p[u]\le\frac{31}{46080}\wo(u)-D_p(u/\alpha),
\end{equation}
where the full nonnegative slack is
\[
 D_p(g)=\sum_{n\ge2}a_n\int_{\{x,x+\ell_n\in O\}}
  \alpha(x)\alpha(x+\ell_n)|g(x+\ell_n)-g(x)|^2\,dx.
\]
Every prime power and both row orientations are retained.
Combining \eqref{eq:inversecost} and \eqref{eq:signedcapmain},
\begin{equation}\label{eq:Slower}
 S[u]\ge\tg[u]+D_p(u/\alpha)-\varepsilon_8\wo(u),\qquad
 \varepsilon_8=\frac{31}{46080}+\frac{13}{5000}
 =\frac{18851}{5760000}.
\end{equation}

For $f\in\D$ use the exterior-mean gauge
\[
 \beta_O=\mu(O)^{-1}\int_Of\,d\mu,\qquad
 u=\kappa\one_O(f-\beta_O),\quad
 v=\kappa\one_I(f-\beta_O).
\]
Appendix~\ref{app:cuts} puts $u$ in $\D_O^0$ and $v$ in $\D_I$.
Moreover
\[
 \wo(u)=\int_O|f-\beta_O|^2\,d\mu
 \le\int_O|f-\mu(f)|^2\,d\mu\le\Var_\mu(f).
\]
The full bridge and local square completion give
\begin{equation}\label{eq:squaremain}
 \E[f]-\tfrac12\Var_\mu(f)=q_I[v-L^{-1}Ju]+S[u].
\end{equation}
The local square is nonnegative by \eqref{eq:Lcoercive}.
Dropping only that square, $\tg$ and $D_p$ in \eqref{eq:Slower} proves
Theorem~\ref{thm:conditionalglobal}.
The exact comparison is
\[
 \frac{2861149}{5760000}-\frac{71501077}{144000000}
 =\frac{24}{125000}>0.
\]
The full-domain claim uses the proved equality of minimum and maximal energy domains, not a change of realization.

\section{The unresolved regularized sign}\label{sec:limits}
Let $D_2$ be the actual $n=2$ summand of $D_p$ and put
$A_2=\tg+D_2$ on the original centered exterior domain.
Equation~\eqref{eq:Slower} implies only
\[
 A_2[u]\le S[u]+\varepsilon_8\wo(u).
\]
It does not imply $S[u]+A_2[u]\ge0$.
In particular the source estimate is in $\wo$, not an $A_2$-relative inverse-response estimate.

The stronger question of nonnegativity on a faithful $A_2^{1/2}$ completion would additionally require that completion and the corresponding continuous extensions to be justified.
No such whole-completion sign, or even a sign on all of the original domain, is proved here.
The full slack $D_p$ is used only on its stated finite-$\wo$ form domain; it is not extended to a weaker topology by assertion.

If $S+A_2\ge0$ were independently proved on the original domain (in particular if proved on a valid larger completion), combining it with the displayed inequality would yield
$S\ge-(\varepsilon_8/2)\wo$ and the coefficient
$5741149/11520000$.
That number has an unproved premise and is not an achieved bound.
A rank-seven lift, if used in such a future sign analysis, would be
$H_7=L+7E_7E_7^*$, not the rank-64 lift $H_{64}=L+9EE^*$ used here.
Neither finite-dimensional response data nor a single positive-loss estimate supplies a continuation to zero loss.
The global half bound and the Riemann hypothesis remain unproved.

\appendix
\section{Elementary full-series bounds}\label{app:scalar}
For $r>0$ set $S_m(r)=\sum_{k=0}^m r^k/k!$.
The positive exponential series and its geometric tail give
\[
 S_m(r)<e^r\le S_m(r)+\frac{r^{m+1}}{(m+1)!}
                  \frac1{1-r/(m+2)}\qquad(0<r<m+2).
\]
The positive logarithm identity is
\begin{equation}\label{eq:atanh}
 \log\frac{1+z}{1-z}=2\sum_{k\ge0}\frac{z^{2k+1}}{2k+1},
 \qquad0<z<1.
\end{equation}
Omitted denominators may be bounded below and powers summed geometrically.
For example
\begin{gather*}
 693/1000<842/1215<\log2<25/36<7/10,\\
 1098/1000<7379/6720<\log3<1+1/12+1/60=11/10,\\
 1609/1000<\log5<13724/8505<81/50.
\end{gather*}
The lower bound for $\log5$ follows from the first eight terms at $z=2/3$,
which exceed respectively
$1.33333,.19753,.05267,.01672,.00578,.00210,.00079,.00030$.
These are exact finite decimals summing to $1.60922$.
The upper bound uses the first three terms and their geometric remainder.
Squaring proves
$\sqrt2>140/99$, $\sqrt3>55/32$, $\sqrt5>324/145$;
these imply \eqref{eq:coefficientcaps}.
Also $\log n/\sqrt n\le2/e<1$, so $a_n\le1$.
The inequality $\Lambda(n)\le n/2$ follows from $\log n<n/2$ for $n\ge2$.

The arctangent addition formula and alternating series give
\[
 \pi>4\sum_{d\in\{2,3\}}\sum_{k=0}^3
       \frac{(-1)^k}{(2k+1)d^{2k+1}}
 =\frac{1538665}{489888}>\frac{157}{50},
\]
where the last difference is $10417/12247200$.
Also
$\frac{22}{7}-\pi=\int_0^1z^4(1-z)^4/(1+z^2)\,dz>0$.
Thus $\pi^2<10$.
Useful Taylor consequences are
\[
 e^{2/5}<3/2,\quad e^{1/5}<5/4,\quad
 e^{8/5}<5,\quad e^3>401/20>20,\quad e^2>73/10.
\]
The bound at $8/5$ follows from the displayed tail formula with $m=6$.

\subsection{The exterior theta envelope}
For $g(t)=(4t^2-6t)e^{-t}$ and $t\ge3$,
\begin{equation}\label{eq:thetatotal}
 \sum_{j\ge2}g(j^2t)<128t^2e^{-4t}.
\end{equation}
Indeed bound each term by $4t^2j^4e^{-tj^2}$.
Consecutive ratios for $j\ge2$ are at most
$(3/2)^4e^{-5t}<1/2$, so twice the first majorant pays the complete tail.
Both $g(t)$ and $t^2e^{-4t}$ decrease for $t>3$.

At $B=4/5$, $S_7(8/5)>99/20$ gives
$t_B>T_0=15543/1000>31/2$.
The logarithm bounds imply $e^B<\sqrt5<9/4$, so $1+e^{-B}>13/9$.
The degree-nine exponential sum at $3$ is
$7280604/362880>401/20$; the degree-four sum at $543/1000$ is
$41297535960801/24000000000000>43/25$.
Hence $e^{15}>20^5(1+5/400)=3240000$ and $e^{T_0}>5572800$.
Since $4T_0^2-6T_0<8731/10$, the leading part of $A(B)$ is below
$78579/724464000$.
The entire tail is below $32768/20^{20}<10^{-12}$, while
\[
 \frac1{9216}-\frac{78579}{724464000}
 =\frac{279936}{9216\cdot724464000}>10^{-8}.
\]
This proves $A<1/9216$.
The logarithmic derivatives of its two summands are at most
\[
 \frac1{1+e^x}+\frac{2(8t-6)}{4t-6}-2t\le6-2t<-25,
 \qquad 4-8t<-25
\]
on $x\ge B$.
Therefore $W(B+h)\le A(B+h)\le Ae^{-25h}$ and
$\int_B^\infty W\le A/25$.

\subsection{A global theta cap and the exterior probability}
Termwise differentiation is justified by locally uniform summable majorants, and gives
\[
 \kappa'(x)=e^{x/2}\sum_{j\ge1}t_jP(t_j)e^{-t_j},
 \qquad P(t)=-8t^2+30t-15.
\]
At zero the first theta term is below $46/49$:
$4\pi^2-6\pi<4(22/7)^2-6(22/7)$, while
$e^\pi>e^{31/10}>22$.
The remaining terms are bounded by $40j^4e^{-3j^2}$,
whose first term is below $1/250$ and successive ratio below $1/10000$.
Thus $\kappa(0)<46/49+40/9999$.
The polynomial $P$ decreases on $[3,\infty)$, with
$P(31/10)=28/25$ and $P(16/5)=-23/25$.
All $j\ge2$ derivative terms are negative.
A positive derivative is therefore possible only on
$0<x<\frac12\log(32/31)<1/62$.
There,
\[
 \kappa'_+(x)<\frac{60}{59}\frac{16}{5}\frac{28}{25}\frac1{22}
 =\frac{2688}{16225}<\frac16.
\]
Hence
\[
 0<\kappa(x)<\frac{46}{49}+\frac{40}{9999}+\frac1{372}
 =\frac{57440653}{60753924}<\frac{19}{20}.
\]
Evenness gives the global assertion and $\sup W<1/2$.
The normalization proved in Appendix~\ref{sec:energy} gives
\[
 \int_0^\infty W
 =\int_0^\infty\frac{\kappa}{2c}\le\frac12\int_0^\infty c\kappa=\frac18.
\]

For completeness, the endpoint cap $\kappa(B)<1/4000$ follows from
$t_B>31/2$, $e^{B/2}<3/2$, $e^{31/2}>5300000$:
\[
 \kappa(B)<\frac{1302}{5300000}
 +\frac{23064}{5300000^4}\frac{10000}{9999}
 <\frac{1302}{5300000}+\frac1{10^9}<\frac1{4000}.
\]
The exponential lower bound is supplied by $S_{24}(31/2)$;
its excess over $5300000$ is
\[
 \frac{37076883949467729399970752217216753}
 {10409396852733332453861621760000}>0.
\]
The tail ratio in this estimate is below $1/10000$.
Furthermore
\[
 -\kappa'-25\kappa
 =e^{x/2}\sum_{j\ge1}t_j(8t_j^2-130t_j+165)e^{-t_j}>0
 \quad(x\ge B),
\]
since at $t=31/2+u$ the polynomial is $8u^2+118u+72$.
It follows that
\[
 \int_B^\infty e^{x/2}\kappa(x)\,dx
 <\frac{(3/2)(1/4000)}{49/2}=\frac3{196000},
 \qquad \mu(O)<\frac3{49000}.
\]

\subsection{All needed interior point estimates}
Since $t(1/10)>19/5>16/5$, the differentiated polynomial proves that
$\kappa$ and $W$ decrease to the right of $1/10$.
Every tail below is the complete bound \eqref{eq:thetatotal}, with the factor $e^{x/2}$ when estimating $\kappa$.

At $x=1/10$,
$t>(157/50)(61/50)>19/5$,
$1+e^{-x}>19/10$, and $e^{19/5}>20(11/5)=44$.
The leading bound for $W$ is $437/1045$; its tail is below
$46208/96000000<1/1000$.
The margin $21/50-437/1045=19/10450$ pays this tail.

At $x=1/4$, $t>(157/50)(79/48)>103/20$,
$e^{1/4}<9/7$ gives $1+e^{-x}>16/9$, and
$e^{103/20}>20(73/10)(29/25)=4234/25$.
The leading bound is $67671/270976$, the full tail is below
$3456/(20^6\cdot7)<1/100000$, and the margin to $1/4$ is
$73/270976$.

At $x=9/20$,
$(157/50)S_3(9/10)=761921/100000>15/2$,
$1+e^{-9/20}>31/20$, and $e^{15/2}>20^2\cdot4=1600$.
The leading bound is $9/124$, the tail is below $7200/20^{10}<1/1000$,
and the margin to $3/40$ is $3/1240$.

For $\kappa(2/5)$, use $t>69/10$, $e^{1/5}<5/4$,
$e^{69/10}>900$.
The leading bound is $207/1000$, and the tail is below
$\frac{19044}{5\cdot900^4}\frac{10000}{9999}<10^{-6}$,
proving $\kappa(2/5)<1/4$.
The first inequality follows from $(157/50)S_3(4/5)>69/10$.

At $x=1/2$ take $t>17/2$, since $(157/50)S_5(1)>17/2$.
Use $e^{1/4}<9/7<4/3$ and $e^{17/2}>20^2\cdot7\cdot(79/48)$.
The leading bound is $952/13825$ and the tail is below
$36992/(3\cdot20^{11})<1/1000$.
The margin $7/100-952/13825=63/55300$ pays the tail.

For $\kappa(3/5)$,
$S_6(6/5)-520/157=88083/12265625>0$ and
$S_5(7/5)-4=1979/46875>0$ give $t>52/5$ and
$e^{52/5}>32000$; also $e^{3/10}<10/7$.
The leading bound is $(10/7)(9256/25)/32000$.
The full tail is below $19840/20^{13}$, by the majorants
$620j^4e^{-10j^2}$ and their ratio below $1/2$.
Their sum is below $1/50$.

At $x=7/10$, take $t>25/2$, using $e^{7/5}>4$.
Also $e^{7/20}<47/33<10/7$ and
$e^{25/2}>20^4(79/48)$.
The leading bound is $33/11060$ and the tail is below
$200000/(7\cdot20^{16})<10^{-6}$.
The margin to $3/1000$ is $9/553000$.

At $x=3/4$,
$(157/50)S_5(3/2)=896627/64000>14$,
$e^{3/8}<3/2$, and $e^{14}>20^4(73/10)=1168000$.
The leading bound is $21/23360$ and the tail is below
$37632/20^{18}<10^{-5}$.
The margin to $1/1000$ is $59/584000$.
All exponential comparisons above follow from the stated positive sums or their geometric remainder formula; no evaluated theta value enters.

\subsection{Gamma constants and all overlap attenuations}
Since $\rho(s)=e^{s/2}/(2\sinh s)$, the inequality
$\rho(s)\le1/(2s)+1/4$ is equivalent to
$2se^{s/2}\le(2+s)\sinh s$.
For odd degree $n$ coefficient comparison reduces to $n\le2^{n-1}$;
for even $n\ge2$ it reduces to $2^{2-n}\le1$.
Thus the power-series difference is nonnegative.
The lower bound $\rho(B)>5/6$ uses $e^{2/5}<3/2$, $e^{8/5}<5$;
the upper bound $\rho(B)<9/10$ uses $e^{2/5}>10/7$, $e^{8/5}>9/2$.

The harmonic error formula gives
\[
 H_n-\log n-\gamma_E
 =\sum_{k=n}^\infty\int_k^{k+1}\left(\frac1x-\frac1{k+1}\right)\,dx
 >\frac1{2(n+1)}.
\]
At $n=5$ and $\log5>8/5$ this yields $\gamma_E<3/5$.
Together with $\log\pi<6/5$, $\pi/2<11/7$ and $\log2<7/10$,
\[
 \cg<\frac65+\frac35+\frac{11}{7}+\frac{21}{10}
 =\frac{383}{70}<\frac{11}{2}.
\]

For the six overlap caps let $\sigma=25/2$.
The reflection exponents exceed $19/8,47/8,35/4$, from the stated logarithm lower bounds.
The positive exponential sums give
\begin{align*}
 e^{19/8}&>7(185/128)>10,\\
 e^{47/8}&>20\cdot7(289/128)>300,\\
 e^{35/4}&>400(73/10)(21/10)=6132>6000.
\end{align*}
For the translations, use
$\sqrt{3/4}<7/8$, $\sqrt{3/5}<4/5$, $\sqrt{4/5}<9/10$, and
\[
 28\cdot3^{12}<4^{12},\qquad
 2000\cdot3^{12}<5^{13},\qquad
 27\cdot4^{12}<2\cdot5^{12}.
\]
These give $(3/4)^\sigma<1/32$, $(3/5)^\sigma<1/500$ and
$(4/5)^\sigma<1/15$, paying every listed pair.

\section{Coercivity of the rank 64 lift}\label{sec:lift}
Let
\[
 e_0=(2B)^{-1/2},\qquad e_n=B^{-1/2}\cos(k_nx),\quad
 k_n=n\pi/B\quad(n\ge1),
\]
be the physical even Neumann basis. Let $P_{64}$ project onto
$e_0,\ldots,e_{63}$ and $H_{64}=L+9P_{64}$. Write $H=H_{64}$ in this appendix. We prove the lift estimate
without using the stored response data.

\subsection{The entire Gamma form dominates a diagonal form}
For $a>0$ the restricted whole-line Green kernel
$e^{-a\abs{x-y}}/(2a)$ is the inverse $R_{R,a}$ of
$-\partial_x^2+a^2$ with Robin form boundary contribution
$a(\abs{v(-B)}^2+\abs{v(B)}^2)$. Removing this nonnegative term gives
the Neumann inverse $R_{N,a}$, so $R_{R,a}\preceq R_{N,a}$.
For every zero-extended local vector,
\[
 \int_0^\infty e^{-as}\norm{v(\cdot+s)-v}_2^2\dd s
 =\frac2a\norm v_2^2-2a\ip v{R_{R,a}v}
 \ge\frac2a\norm v_2^2-2a\ip v{R_{N,a}v}.
\]
Using $\rho(s)=\sum_{j\ge0}e^{-\lambda_js}$, where
$\lambda_j=2j+1/2$, and monotone summation of nonnegative forms, yields
\begin{equation}\label{eq:DN}
 t_\Gamma\succeq D_N,\qquad
 D_Ne_n=d_ne_n,\qquad
 d_n=2\sum_{j\ge0}\frac{k_n^2}{\lambda_j(\lambda_j^2+k_n^2)}.
\end{equation}
This is a form inequality on the whole local minimum domain. It does
not impose Neumann boundary conditions on an original vector. Lower
semicontinuity extends the core comparison, with inclusion in the
form domain of $D_N$.

The numbers $d_n$ are nonnegative and increasing. We give a short
rational proof that $d_{64}>9$. Set
\[
 f_t(a)=\frac{2t^2}{a(a^2+t^2)},\qquad G(t)=\sum_{j\ge0}f_t(2j+1/2).
\]
Each summand increases in $t$, and $k_{64}=80\pi>240$.
Differentiation gives
\[
 f_t''(a)=\frac{4t^2(6a^4+3a^2t^2+t^4)}{a^3(a^2+t^2)^3}>0.
\]
The mesh-two trapezoidal inequality for this convex decreasing function
therefore gives
\begin{align*}
 G(240)\ge{}&f_{240}(1/2)+f_{240}(5/2)+f_{240}(9/2)
 +\tfrac12f_{240}(13/2)\\
 &+\tfrac12\log\bigl(1+(480/13)^2\bigr),
\end{align*}
since $\int_A^\infty f_t(a)\dd a=\log(1+t^2/A^2)$.
The inequality $(1+x)^{-1}>1-x$ proves that the four rational terms
exceed, respectively,
\[
 \frac{3999}{1000},\qquad\frac{799}{1000},\qquad
 \frac{444}{1000},\qquad\frac{153}{1000}.
\]
For example they exceed $4-1/57600$, $4/5-5/57600$,
$4/9-1/6400$ and $2/13-13/115200$, respectively.
Also \eqref{eq:atanh} gives
\[
 \log2>\frac{842}{1215}>\frac{693}{1000},\qquad
 \log(15/13)>\frac17.
\]
It follows that
\[
 \tfrac12\log\bigl(1+(480/13)^2\bigr)
 >\log(480/13)=5\log2+\log(15/13)
 >\frac{693}{200}+\frac17.
\]
Consequently
\begin{equation}\label{eq:d64}
 d_{64}>G(240)>\frac{3151}{350}=9+\frac1{350}>9.
\end{equation}

\subsection{A whole operator bound for the local prime term}
Write $s=\log2$, $t=\log3$, $d=B-s$ and
$a=s/\sqrt2$, $b=s/2$, $a_3=t/\sqrt3$.
The elementary log bounds give $.69<s<.70$, $1.09<t<1.10$ and
$\log5>2B$. The lower bound for $s$ follows from $842/1215>.693$;
for $t$, the first three terms in \eqref{eq:atanh} give
$\log3>1+1/12+1/80>1.09$.
Under the unitary fold on $(0,B)$ partition that interval
(up to endpoints) into
\[
 C=(0,d),\quad L_-=(s-d,s),\quad L_+=(s,s+d),\quad
 J_0=(d,s-d).
\]
Prime 2 connects $y\in C$ to $s-y$ and $s+y$, with weight $a$,
and reflects $J_0$. Prime 4 connects the two leaves $L_-,L_+$
with weight $b$. Prime 3 reflects $S=(t-B,B)$ with weight $a_3$.
The inequalities
$2s>t$, $t+s>2B$, $t+2B<4s$ show that the leaves lie in $S$,
that their prime-3 images lie in $J_0$, and that $C$ does not meet $S$.

Define the positive number
\[
 R_p=\frac{b+a_3+\sqrt{(b+a_3)^2+8a^2}}2,
 \qquad\eta_0=2a/R_p,
 \qquad w_0=\eta_0 1_C+1_{J_0\cup L_-\cup L_+}.
\]
The weighted row ratio $(\PP_Iw_0)/w_0$ equals $R_p$ on the center
and leaves and is at most $a+a_3$ on $J_0$. Here
$\PP_I$ is the complete local prime operator in the definition of $L$.
One has $R_p>5/4$: its defining quadratic at $5/4$ is negative using
$b>1/3$, $a_3>4/7$, $2a^2>4/9$.
Moreover $a<1/2$, $a_3<121/190$ and hence $a+a_3<5/4$.
Every edge map is measure-preserving and every edge weight is
nonnegative. Applying
$2\abs{uv}\le(w_y/w_x)u^2+(w_x/w_y)v^2$ on each symmetric edge,
and summing, proves
$\abs{\ip v{\PP_Iv}}\le R_p\norm v_2^2$.
This retains all low/high interactions of the full operator.
Finally
\[
 b+a_3<\frac7{20}+\frac{121}{190}=\frac{75}{76},\quad
 2a^2<\frac{49}{100},\quad
 \left(\frac{27}{20}\right)^2-\frac{75}{76}\frac{27}{20}
 -\frac{49}{100}=\frac1{3800}>0.
\]
Thus
\begin{equation}\label{eq:primeop}
 \norm{\PP_I}<\frac{27}{20}.
\end{equation}

\subsection{The lift constant and a residual majorant}
The elementary estimate
\begin{equation}\label{eq:cgamma}
 c_\Gamma<\frac65+\frac35+\frac{11}{7}+\frac{21}{10}
 =\frac{383}{70}
\end{equation}
uses $e^{6/5}>401/125>22/7$ and $\gamma_E<3/5$.
To check the latter without evaluating Euler's constant, the harmonic
error identity gives
\[
 H_n-\log n-\gamma_E
 =\sum_{k=n}^\infty\int_k^{k+1}
       \left(\frac1x-\frac1{k+1}\right)\dd x
 >\frac1{2(n+1)}.
\]
Take $n=5$ and $\log5>8/5$ to get
$\gamma_E<H_5-8/5-1/12=3/5$.

\begin{theorem}[Rank-64 lift coercivity]\label{thm:lift}
On the complete mean-inclusive local minimum domain,
\begin{equation}\label{eq:lift}
 H_{64}=L+9P_{64}\succeq\frac{61}{28}I,
 \qquad H^{-1}\preceq\frac{28}{61}I.
\end{equation}
This is a coercivity statement for $H$, not for the unlifted $L$.
\end{theorem}
\begin{proof}
Dropping only the nonnegative pole in a lower estimate, equations
\eqref{eq:DN} and \eqref{eq:primeop} give
\[
 L\succeq D_N-(c_\Gamma+27/20)I.
\]
Below mode 64 the lifted diagonal is at least $9-c_\Gamma-27/20$;
from mode 64 onward the same bound follows from \eqref{eq:d64}.
Since $9-383/70-27/20=61/28$, \eqref{eq:lift} follows.
The pole has not been removed from the operator being inverted.
\end{proof}
For the residual contract define the larger diagonal majorant
\begin{equation}\label{eq:Dold}
 D_{\rm old}e_n=\beta_ne_n,\qquad
 \beta_n=\begin{cases}(d_n+1/2)^{-1},&n<64,\\
 (d_n-17/2)^{-1},&n\ge64.
 \end{cases}
\end{equation}
Translation Cauchy--Schwarz gives
$\norm{\PP_I}\le2(a_2+a_3+a_4)<3$, and
$c_\Gamma<11/2$, so $H\succeq D_N-17I/2+9P_{64}$.
Equation~\eqref{eq:d64} proves
$0<H^{-1}\preceq D_{\rm old}\preceq2I$.
Contract (a) uses this larger majorant. The sharper scalar bound for $H^{-1}$ is used separately in the inverse-cost estimate.

\section{An exact signed exterior prime identity}\label{sec:signed}
In this appendix write $A_B=A$.
This section is analytic and requires no stored numerical input or
exterior moment constraint. Set $g=u/\alpha$ and
$E_n=\{x:x,x+\ell_n\in O\}$. Define the nonnegative full-prime slack
\begin{equation}\label{eq:slack}
 D_p(g)=\sum_{n\ge2}a_n\int_{E_n}
 \alpha(x)\alpha(x+\ell_n)[g(x+\ell_n)-g(x)]^2\dd x.
\end{equation}
The sum includes all prime powers. The two row orientations arise from
the two diagonal terms in each edge square, as the proof below shows.

\subsection{The complete exterior row and its one reflection}
On the positive exterior fold $f=\sqrt2u$, the full prime expression is
\begin{equation}\label{eq:exteriorfold}
 (\PP_Of)(x)=\sum_{n\ge2}a_n\big[
 f(x+\ell_n)+1_{x>B+\ell_n}f(x-\ell_n)
 +1_{B<x<\ell_n-B}f(\ell_n-x)\big],\quad x>B.
\end{equation}
The first two terms are same-side orientations. The third is one
reflected opposite-side input, with one coefficient $a_n$. The cases
for $x-\ell_n$ are disjoint: positive exterior, killed interval and
negative exterior. Since $\log4<2B<\log5$, no reflection occurs for
$n=2,3,4$. The unitary fold already absorbs the factor two from the
two half-lines.

The relation $p(u,u)=\int_B^\infty f\PP_Of$ is understood as an
absolutely convergent form or weighted-duality pairing, first for
truncated sums. It is not an assertion that $\PP_Of$ belongs to
physical $L^2$ for every finite-$\wo$ vector. Equivalently it is the
$L^2$ pairing with the bounded conjugate
$M_\alpha\PP_OM_\alpha$, justified by the following row bound.
Define the nonnegative potential on all of $O$ by
\begin{equation}\label{eq:primepotential}
 V_p(x)=\alpha(x)\sum_{n\ge2}a_n
 \big[\alpha(x+\ell_n)1_O(x+\ell_n)
      +\alpha(x-\ell_n)1_O(x-\ell_n)\big].
\end{equation}
For $x>B$, either existing same-side product is at most
$A_Bn^{-25/2}$: the sum of the two distances from $B$ is at least
$\ell_n$. On the reflected branch it is at most
$A_Be^{-(25/2)(\ell_n-2B)}$, because those distances sum exactly to
$\ell_n-2B$. There is only one such branch in a row.

The two same-side budgets together satisfy
\[
 2\sum_{n\ge2}a_nn^{-25/2}\le2\sum_{n\ge2}n^{-12}
 \le2\left(2^{-12}+\int_2^\infty t^{-12}\dd t\right)
 =\frac{13}{22528}<\frac1{1000}.
\]
At $n=5$, equations \eqref{eq:coefficientcaps} and
$e^{-(25/2)(\log5-2B)}<9/10$ give cost less than $261/400$.
There is no $n=6$ coefficient. Every remaining reflection is covered by
\eqref{eq:bothtails}, since
$a_ne^{-(25/2)(\log n-2B)}\le(5/n)^{12}$ for $n\ge7$.
Consequently, for all $x\in O$,
\begin{equation}\label{eq:primecap}
 V_p(x)<\left(\frac1{1000}+\frac{261}{400}+\frac3{100}\right)A_B
 =\frac{1367}{2000}A_B<\frac7{10}A_B.
\end{equation}
Evenness extends the pointwise bound to the negative half-line without
doubling it. The summable majorants give uniform row convergence and
boundedness of the conjugated operator, including every prime tail.

\subsection{The signed identity and the analytic cap}
\begin{theorem}[Full-prime slack and mass bound]\label{thm:signed}
For every real even exterior $u$ of finite $\wo$, one has
\begin{equation}\label{eq:slackidentity}
 p(u,u)+D_p(u/\alpha)
 =\int_O V_p(x)|u(x)/\alpha(x)|^2\dd x,
\end{equation}
with $0\le D_p(u/\alpha)\le(7/5)A_B\wo(u)$.
Put
\[
 V(x)=c_\Gamma\alpha(x)^2+V_p(x),\qquad N=\frac{31}{46080}.
\]
Then $0<V(x)<N$ and
\begin{equation}\label{eq:negativeanalytic}
 c_\Gamma\norm u_2^2+p(u,u)
 \le N \wo(u)-D_p(u/\alpha).
\end{equation}
Whenever $t_\Gamma(u)$ is finite, the exact spatial identity is
\begin{equation}\label{eq:signedq0}
 q_0(u)=t_\Gamma(u)+D_p(u/\alpha)
              -\int_OV(x)|u(x)/\alpha(x)|^2\dd x.
\end{equation}
The cap is an analytic conclusion, not a numerical premise.
\end{theorem}
\begin{proof}
For a finite sum, expansion of one edge square in \eqref{eq:slack}
gives cross term $-2a_nu(x)u(x+\ell_n)$. The two diagonal terms,
after a change of variable in one, give the two orientations in
\eqref{eq:primepotential}. This proves \eqref{eq:slackidentity} for
truncated sums with exactly the factor in the definition of $p$.
The inequality $(a-b)^2\le2(a^2+b^2)$ gives
\[
 0\le D_p(g)\le2\int_OV_pg^2\le\frac75 A_B\norm g_2^2.
\]
The same row bound controls the absolute prime bilinear form by
weighted Cauchy--Schwarz. Monotone convergence of the slack series and
absolute convergence of the prime pairing extend the identity to all
$g\in L^2(O)$. Thus negative prime correlation increases the retained
slack; its sign has not been replaced by an absolute correlation.

Equation~\eqref{eq:cgamma} gives $c_\Gamma<11/2$. Together with
\eqref{eq:envelope} and \eqref{eq:primecap},
\[
 0<V(x)<\left(\frac{11}{2}+\frac7{10}\right)A_B
 =\frac{31}{5}A_B<\frac{31}{46080}.
\]
Adding the mass term proves \eqref{eq:negativeanalytic}, and
subtracting the exact identity from $t_\Gamma$ proves
\eqref{eq:signedq0}. The slack and spatial potential are $\wo$-bounded,
so the same identities extend in the original form topology, in
particular on the fixed minimum domain $\D_O^0$.
\end{proof}

The slack has trivial kernel on the even weighted exterior space.
Indeed its $n=2$ same-side edges have strictly positive weight, so
zero slack forces $g(x+\log2)=g(x)$ almost everywhere on $(B,\infty)$.
The $L^2$ norm is then an infinite sum of equal nonnegative period
integrals and can be finite only if each vanishes. Evenness supplies
the other half-line. This trivial-kernel statement does not assert a
uniform inequality $D_p\ge\delta\wo$ with $\delta>0$.


\section{The theta kernel and the original closed energy}\label{sec:energy}
In the foundation proofs write $s_n=\ell_n=\log n$.

For $x\ge0$, define
\begin{equation}\label{eq:kappa}
 \kappa(x)=e^{x/2}\sum_{j\ge1}(4\tau_j(x)^2-6\tau_j(x))e^{-\tau_j(x)},
 \qquad \tau_j(x)=\pi j^2e^{2x},
\end{equation}
and extend evenly. We record the normalization and decay needed later.

\begin{lemma}[Smoothness and normalization]\label{lem:theta}
The kernel $\kappa$ is smooth, strictly positive, and even. There is a finite $C_\kappa>0$ such that
\begin{equation}\label{eq:envelopek}
 \kappa(x)\le C_\kappa\exp\!\left(\frac92|x|-\pi e^{2|x|}\right),
 \qquad
 C_\kappa=4\pi^2\sum_{j\ge1}j^4e^{-\pi(j^2-1)}.
\end{equation}
Furthermore $\int_\R2c\kappa\,dx=1$, and
\begin{equation}\label{eq:envelopeW}
 0<W(x)\le C_\kappa\exp(4|x|-\pi e^{2|x|}).
\end{equation}
\end{lemma}
\begin{proof}
The series in~\eqref{eq:kappa}, together with every derivative, converges uniformly on compact subsets of $[0,\infty)$. Each summand is positive because $\tau_j\ge\pi>3/2$. Dropping the negative polynomial term and using $e^{2x}\ge1$ gives~\eqref{eq:envelopek} for $x\ge0$.

To justify smooth evenness at zero, rather than assume it from reflection, set
\[
 H(x)=e^{x/2}\sum_{j\ge1}e^{-\pi j^2e^{2x}},\qquad x\in\R.
\]
The Jacobi transformation \cite[Eq.~20.7.32]{DLMF} implies
$H(-x)=H(x)+\sinh(x/2)$. Applying $D^2-1/4$ shows that $H''-H/4$ is even; direct differentiation identifies it with~\eqref{eq:kappa} on the positive half-line. Thus it gives the claimed smooth extension. Differentiating the transformation at zero yields $H'(0)=-1/4$. Since $c''=c/4$, integration by parts and the rapid decay of $H,H'$ at $+\infty$ give
\[
 \int_0^\infty c(x)\kappa(x)\,dx
  =[cH'-c'H]_0^\infty=\frac14.
\]
Evenness proves the normalization. Finally $2c(x)\ge e^{|x|/2}$ gives~\eqref{eq:envelopeW}.
\end{proof}

Put $\rho(s)=e^{-s/2}/(1-e^{-2s})$ for $s>0$, and $a_n=\Lambda(n)/\sqrt n$. The full positive energy is
\begin{align}
 \E_\Gamma[f]&=\int_0^\infty\rho(s)\int_\R
       \kappa(x)\kappa(x+s)|f(x+s)-f(x)|^2\,dx\,ds,\label{eq:Eg}\\
 \E_P[f]&=\sum_{n\ge2}a_n\int_\R
       \kappa(x)\kappa(x+s_n)|f(x+s_n)-f(x)|^2\,dx,\label{eq:Ep}\\
 \E[f]&=\E_\Gamma[f]+\E_P[f],\qquad
 \D=\{f\in\He:\E[f]<\infty\}.\label{eq:Edomain}
\end{align}
These are nonnegative integrals and sums, interpreted by Tonelli's theorem; no cancellation defines their domain. Write
\[
 \mu(f)=\int_\R f\,d\mu,\qquad
 \Var_\mu(f)=\norm{f}_\mu^2-|\mu(f)|^2.
\]

\begin{proposition}[Maximal domain and compact smooth form core]\label{prop:core}
The form $\E$ on $\He$ is densely defined, nonnegative, and closed. Its maximal integral domain $\D$ is the closure of even $C_c^\infty(\R)$ in the norm $(\norm{f}_\mu^2+\E[f])^{1/2}$. Constants belong to $\D$, and $\E[f]=0$ holds precisely for almost-everywhere constant $f$.
\end{proposition}
\begin{proof}
The theta envelope implies that for every $A_0>0$ there is a constant $C_{A_0}$ such that $\kappa(x)\le C_{A_0}e^{-A_0|x|}$. Consequently
\begin{equation}\label{eq:edgefinite}
 \int_\R\kappa(x)\kappa(x+s)\,dx
 \le C_{A_0}^2(s+A_0^{-1})e^{-A_0s},\qquad s\ge0.
\end{equation}
Since $a_n\le n$, the complete prime edge measure has finite mass, by taking $A_0>2$ in~\eqref{eq:edgefinite}. For every $L<\infty$,
\begin{equation}\label{eq:rhointegrable}
 \int_0^\infty\rho(s)\min\{1,L^2s^2\}\,ds<\infty.
\end{equation}
These observations admit compact smooth functions and constants to the domain.

For closedness, let $(f_j)$ be Cauchy in the form norm, with $L^2(\mu)$ limit $f$. A subsequence converges almost everywhere for $\mu$, hence for Lebesgue measure, because the density is strictly positive. A countable union of translated null sets is null, so all prime differences converge almost everywhere. Fubini's theorem gives convergence for the Gamma edge measure as well. For each fixed $k$, Fatou's lemma gives
\[
 \E[f-f_k]\le\liminf_{j\to\infty}\E[f_j-f_k].
\]
It follows that $f\in\D$ and $f_k\to f$ in form norm.

We prove the core assertion first for real profiles. Let $T_M(t)=\max\{-M,\min\{t,M\}\}$. The remainder $r_M(t)=t-T_M(t)$ is $1$-Lipschitz, tends pointwise to zero, and satisfies $|r_M(t)|\le|t|$. Thus dominated convergence in both edge measures and in $L^2(\mu)$ gives $T_M(f)\to f$ in form norm. It suffices to approximate bounded even profiles.

Choose smooth even cutoffs $0\le\chi_R\le1$ equal to one on $[-R,R]$, supported in $[-2R,2R]$, and with uniformly bounded Lipschitz constants. Equation~\eqref{eq:rhointegrable}, the finite prime-edge mass, and dominated convergence give $\E[\chi_R]\to0$. For bounded $f$, the identity for a product difference yields
\begin{align*}
 |[f(1-\chi_R)](y)-[f(1-\chi_R)](x)|^2
 &\le2(1-\chi_R(y))^2|f(y)-f(x)|^2\\
 &\quad+2\norm{f}_\infty^2|\chi_R(y)-\chi_R(x)|^2.
\end{align*}
Hence $f\chi_R\to f$ in form norm. It remains to smooth a compactly supported profile $g\in\D$.

Suppose $\supp g\subset[-K,K]$. Its ordinary $L^2$ norm is finite, since the density of $\mu$ is bounded below on this interval. For $0<s\le1$, a nonzero difference of $g$ has both endpoints in $[-K-1,K+1]$. The positive lower bound of $\kappa$ there controls the small-shift part of $\tg[g]$ by $\E_\Gamma[g]$. The part with $s>1$ is bounded by $4\norm{g}_2^2\int_1^\infty\rho(s)\,ds$. Thus $\tg[g]<\infty$.

Let $g_\epsilon=\zeta_\epsilon*g$, where $\zeta_\epsilon$ is an even compactly supported smooth approximate identity. Translations commute with convolution, and
\[
 \norm{\Delta_s(g_\epsilon-g)}_2\le2\norm{\Delta_sg}_2,
 \qquad\Delta_sg=g(\cdot+s)-g.
\]
Dominated convergence gives $\tg[g_\epsilon-g]\to0$. Since $\kappa$ is bounded,
$\E_\Gamma[g_\epsilon-g]\le\norm{\kappa}_\infty^2\tg[g_\epsilon-g]\to0$.
All the smoothed profiles have support in one fixed compact interval $J$. For any $h$ supported there,
\begin{equation}\label{eq:compactprimerate}
 \E_P[h]\le2\int_J\kappa(x)|h(x)|^2
     \sum_{n\ge2}a_n\bigl(\kappa(x-s_n)+\kappa(x+s_n)\bigr)\,dx
 \le C_J\norm{h}_2^2.
\end{equation}
Indeed, the complete rate in parentheses is uniformly bounded on $J$: use $\kappa(x\pm\log n)\le C_{A_0}e^{A_0\sup_J|x|}n^{-A_0}$ and $A_0>2$. Ordinary $L^2$ convergence therefore proves prime-energy convergence. The local upper bound of the density gives $L^2(\mu)$ convergence. Applying the argument to real and imaginary parts proves the complex core assertion.

Finally $\E[1]=0$. If $\E[f]=0$, the strictly positive Gamma density implies $f(y)=f(x)$ for almost every pair $x<y$. Fubini's theorem forces $f$ to be almost everywhere constant. Density of the core in $\He$ is immediate from the positive smooth density of $\mu$.
\end{proof}

\section{The explicit formula and the compact physical bridge}\label{sec:bridge}

We now prove the identity that connects the positive energy to the physical form. The classical explicit formula is used as a theorem, without any assumption about the location of nontrivial zeta zeros inside the critical strip. We first fix its normalization, then justify its use with the noncompact theta kernel, and finally perform the ground-state calculation with all endpoints retained.

In this section use the nonunitary Fourier transform
\[
 \mathcal Fu(z)=\int_\R u(x)e^{izx}\,dx,
 \qquad m(u)=\int_\R c(x)u(x)\,dx.
\]
Let $\rho_\zeta$ range over the distinct nontrivial zeros of $\zeta$, with multiplicity $\nu_{\rho_\zeta}$, and put $w_{\rho_\zeta}=(\rho_\zeta-1/2)/i$. Thus $|\operatorname{Im}w_{\rho_\zeta}|<1/2$; no reality of $w_{\rho_\zeta}$ is asserted. We use the classical zero-counting bound
\begin{equation}\label{eq:zerocount}
 \sum_{|\operatorname{Im}\rho_\zeta|\le T}\nu_{\rho_\zeta}
       =O(T\log(T+2))\qquad(T\ge2).
\end{equation}
This weak form of the classical Riemann--von Mangoldt estimate \cite[Theorem~0.1]{GarrettZeros} is sufficient for every zero sum below.

\begin{theorem}[Classical Guinand explicit formula]\label{thm:guinand}
For $g\in C_c^\infty(\R)$,
\begin{align}\label{eq:guinand}
 \sum_{\rho_\zeta}\nu_{\rho_\zeta}\mathcal Fg(w_{\rho_\zeta})
 &=\mathcal Fg(i/2)+\mathcal Fg(-i/2)\notag\\
 &\quad+\frac1{2\pi}\int_\R
   \bigl(\operatorname{Re}\psi(1/4+it/2)-\log\pi\bigr)
       \mathcal Fg(t)\,dt\notag\\
 &\quad-\sum_{n\ge2}a_n\bigl(g(s_n)+g(-s_n)\bigr).
\end{align}
The zero sum is absolutely convergent. This is the classical formula in the Fourier convention of Garrett \cite[Theorem~0.1]{Garrett}, with prime powers collected by the von Mangoldt function.
\end{theorem}

The statement is quoted, not proposed as a new result. Absolute convergence follows from the rapid decay of $\mathcal Fg$ in the closed strip $|\operatorname{Im}z|\le1/2$ and~\eqref{eq:zerocount}.

\begin{lemma}[Bilinear physical normalization]\label{lem:bilinear}
For real even $u,v\in C_c^\infty(\R)$,
\begin{equation}\label{eq:bilinearzeros}
 q_0(u,v)+2m(u)m(v)
   =\sum_{\rho_\zeta}\nu_{\rho_\zeta}
       \mathcal Fu(w_{\rho_\zeta})\mathcal Fv(w_{\rho_\zeta}).
\end{equation}
\end{lemma}
\begin{proof}
The digamma integral \cite[Eq.~5.9.16]{DLMF} gives, after subtracting its value at $1/4$ and setting its integration variable equal to $2s$,
\begin{equation}\label{eq:digammamultiplier}
 \operatorname{Re}\psi(1/4+it/2)-\psi(1/4)
    =2\int_0^\infty\rho(s)(1-\cos(ts))\,ds=:d_\Gamma(t).
\end{equation}
The difference integral is convergent at zero and infinity. Plancherel and polarization identify
\[
 \tg(u,v)=\frac1{2\pi}\int_\R
       d_\Gamma(t)\mathcal Fu(t)\mathcal Fv(t)\,dt
\]
for real even tests. Apply Theorem~\ref{thm:guinand} to $g=u*v$. Its Fourier transform is the product of the two transforms, and $g(0)=\int uv$. Evenness gives
\[
 \mathcal Fu(i/2)=\mathcal Fu(-i/2)=m(u),\qquad
 g(s_n)+g(-s_n)
  =\int_\R\bigl(u(x+s_n)v(x)+u(x)v(x+s_n)\bigr)\,dx.
\]
These are respectively the pole and the polarized prime term. The constant in the Gamma term is $\psi(1/4)-\log\pi=-\cg$, proving~\eqref{eq:bilinearzeros}.
\end{proof}

For complex even tests the zero-side product is
$\mathcal Fu(w)\overline{\mathcal Fv(\overline w)}$. It is not generally $|\mathcal Fu(w)|^2$ when $u=v$ and $w$ is nonreal. We need only the real bilinear identity and its conjugate-linear extension in the second argument; no positivity is inferred from this zero sum.

\begin{lemma}[Entire theta transform]\label{lem:thetaxi}
Let
\[
 \xi(s)=\tfrac12s(s-1)\pi^{-s/2}\Gamma(s/2)\zeta(s).
\]
Then the entire continuation of this function satisfies
\begin{equation}\label{eq:thetaxi}
 \mathcal F\kappa(z)=\xi(1/2+iz),\qquad z\in\C.
\end{equation}
In particular, $\mathcal F\kappa(w_{\rho_\zeta})=0$ for every nontrivial zero, whether or not it is on the critical line.
\end{lemma}
\begin{proof}
Write $\Theta(t)=\sum_{j\ge1}e^{-\pi j^2t}$, so that the function in Lemma~\ref{lem:theta} is $H(x)=e^{x/2}\Theta(e^{2x})$. For $\operatorname{Re}s>1$, termwise integration gives
\[
 \pi^{-s/2}\Gamma(s/2)\zeta(s)
   =\int_0^\infty\Theta(t)t^{s/2-1}\,dt.
\]
The Jacobi transformation is
$\Theta(t)=t^{-1/2}\Theta(1/t)+(t^{-1/2}-1)/2$.
Split the integral at one and substitute $t\mapsto1/t$ in its lower part. The elementary remainder integrates to $1/[s(s-1)]$, and hence
\begin{align}\label{eq:mellintheta}
 \pi^{-s/2}\Gamma(s/2)\zeta(s)
 &=\frac1{s(s-1)}+I(s),\notag\\
 I(s)&=\int_1^\infty\Theta(t)
       \bigl(t^{s/2-1}+t^{(1-s)/2-1}\bigr)\,dt.
\end{align}
The integral $I$ is entire by local uniform double-exponential domination after $t=e^{2x}$. Thus
$\xi(s)=1/2+s(s-1)I(s)/2$ is entire. For $s=1/2+iz$, the same substitution gives
\begin{equation}\label{eq:xiH}
 \xi(1/2+iz)=\frac12
       -2(z^2+1/4)\int_0^\infty H(x)\cos(zx)\,dx.
\end{equation}
Because $\kappa=H''-H/4$ on the positive half-line, evenness and two integrations by parts give
\[
 \mathcal F\kappa(z)
  =-2H'(0)-2(z^2+1/4)\int_0^\infty H(x)\cos(zx)\,dx.
\]
All boundary terms at infinity vanish for each complex $z$. The value $H'(0)=-1/4$ from Lemma~\ref{lem:theta} proves~\eqref{eq:thetaxi}. The defining prefactors of $\xi$ are finite and nonzero at a nontrivial zero of $\zeta$, so the last assertion follows.
\end{proof}

\begin{proposition}[Weak theta source identity]\label{prop:thetanull}
For every even $h\in C_c^\infty(\R)$,
\begin{equation}\label{eq:thetanull}
 q_0(\kappa,h)=-\overline{m(h)}.
\end{equation}
All terms on the left exist separately as convergent form pairings. Equivalently, $(D_\Gamma-\cg-P)\kappa=-c$ on even compact smooth tests. This is a weak identity and does not assert that $c$ belongs to ordinary $L^2$.
\end{proposition}
\begin{proof}
We first take $h$ real. Let $\chi_R(x)=\chi(x/R)$, $R\ge1$, where $\chi$ is even and smooth, equals one on $[-1,1]$, has support in $[-2,2]$, and satisfies $0\le\chi\le1$. Put $k_R=\kappa\chi_R$. Differentiating the theta series on each half-line shows that $\kappa$ and its derivatives decrease faster than every exponential. Consequently $k_R\to\kappa$ in $H^1(\R)$ and in $L^1(e^{|x|/2}dx)$. For $v\in H^1(\R)$,
\[
 \tg[v]\le\norm{v'}_2^2\int_0^1\rho(s)s^2\,ds
              +4\norm{v}_2^2\int_1^\infty\rho(s)\,ds.
\]
It follows that the Gamma pairing, the ordinary $L^2$ pairing, and the moment in~\eqref{eq:bilinearzeros} converge as $R\to\infty$.

The complete prime pairing also converges without a spatially uniform incoming-rate assumption. If $\supp h\subset[-K,K]$, changing variables in one of its two terms gives
\[
 p(k_R,h)=\sum_{n\ge2}a_n\int_{-K}^K h(x)
           \bigl(k_R(x+s_n)+k_R(x-s_n)\bigr)\,dx.
\]
For any fixed $A_0>2$,
\begin{equation}\label{eq:bridgeprimerate}
 \sup_{|x|\le K}\sum_{n\ge2}a_n
       \bigl(\kappa(x+s_n)+\kappa(x-s_n)\bigr)
 \le2C_{A_0}e^{A_0K}\sum_{n\ge2}a_nn^{-A_0}<\infty.
\end{equation}
Here $a_n\le n$ suffices. Since $|k_R|\le\kappa$, dominated convergence proves the assertion and absolute convergence of the limiting pairing.

For the zero side, uniformly in $R$ and in $|\operatorname{Im}z|\le1/2$,
\[
 |\mathcal Fk_R(z)|\le\int_\R\kappa(x)e^{|x|/2}\,dx<\infty.
\]
Repeated integration by parts in the transform of the fixed compact test gives
$|\mathcal Fh(z)|\le C_{h,N}(1+|\operatorname{Re}z|)^{-N}$ throughout this strip. Equation~\eqref{eq:zerocount} therefore supplies a summable majorant for the products in~\eqref{eq:bilinearzeros}, for example with $N=3$. Pointwise transform convergence and Lemma~\ref{lem:thetaxi} make their limiting sum zero. Thus
\[
 q_0(\kappa,h)+2m(\kappa)m(h)=0.
\]
Lemma~\ref{lem:theta} gives $m(\kappa)=1/2$, proving~\eqref{eq:thetanull} for real $h$. Splitting an arbitrary $h$ into real and imaginary parts and using conjugate linearity in the second argument proves the stated version.
\end{proof}

\begin{theorem}[Compact physical bridge]\label{thm:compactbridge}
For every even $f\in C_c^\infty(\R)$,
\begin{equation}\label{eq:provedbridge}
 q_0[\kappa f]=\E[f]-\frac12\norm{f}_\mu^2.
\end{equation}
\end{theorem}
\begin{proof}
Set $h=\kappa|f|^2$, which is a real even compact smooth test. For $a=\kappa(x)$, $b=\kappa(x+s)$, $F=f(x)$, and $G=f(x+s)$, the pointwise identity
\[
 |bG-aF|^2-(b-a)(b|G|^2-a|F|^2)=ab|G-F|^2
\]
gives, after integration,
\begin{equation}\label{eq:gammaground}
 \tg[\kappa f]-\tg(\kappa,h)=\E_\Gamma[f].
\end{equation}
This is an equality between finite quantities: $\kappa$, $h$, and $\kappa f$ are in $H^1$, their Gamma pairings are finite by the estimate in the preceding proof, and $\E_\Gamma[f]$ is finite by Proposition~\ref{prop:core}.

For the primes the two endpoint terms are both essential. Their full expansion gives
\begin{align}\label{eq:primeground}
 p[\kappa f]-p(\kappa,h)
 &=\sum_{n\ge2}a_n\int_\R\kappa(x)\kappa(x+s_n)
       \bigl(2\operatorname{Re}(f(x+s_n)\overline{f(x)})\notag\\
 &\hspace{44mm}-|f(x+s_n)|^2-|f(x)|^2\bigr)\,dx\notag\\
 &=-\E_P[f].
\end{align}
The compact-profile prime energy is finite, the compact correlation $p[\kappa f]$ has only finitely many nonzero shifts, and $p(\kappa,h)$ is absolutely convergent by~\eqref{eq:bridgeprimerate}. Thus no difference of divergent prime expressions is used. Finally,
$\norm{\kappa f}_2^2=\ip{\kappa}{h}_2$ cancels the constant term. Combining~\eqref{eq:gammaground}, \eqref{eq:primeground}, and~\eqref{eq:thetanull} yields
\[
 q_0[\kappa f]=\E[f]-m(h)
             =\E[f]-\int_\R c\kappa|f|^2\,dx,
\]
which is~\eqref{eq:provedbridge}.
\end{proof}

The theorem uses the classical explicit formula and the theta transform, not a zero-location hypothesis. It establishes an identity of forms, not a positivity statement for $q_0$ or for the centered energy shifted by $1/2$.

\section{The complete weighted prime form}\label{sec:prime}
In this appendix $\norm u_\omega^2=\int_\R |u|^2/W\,dx$ is the full-line weighted numerator norm.

The relevant prime bound is in the numerator's weighted norm, not a uniform vanishing estimate for an incoming jump rate. Let $\beta=\sqrt W$ and let $\tau_sh(x)=h(x+s)$ on ordinary $L^2(\R)$.

\begin{lemma}[Weighted shift estimate]\label{lem:prime}
For a finite constant $K$ and every $s\ge0$,
\begin{equation}\label{eq:betaprod}
 \sup_{x\in\R}\beta(x)\beta(x+s)\le K e^{2s-\pi e^s}.
\end{equation}
The self-adjoint series
\begin{equation}\label{eq:KP}
 K_P=\sum_{n\ge2}a_nM_\beta(\tau_{s_n}+\tau_{-s_n})M_\beta
\end{equation}
converges absolutely in operator norm and preserves even functions. If $K_{P,N}$ denotes its head through $n=N$, then
\begin{equation}\label{eq:rN}
 \norm{K_P-K_{P,N}}\le r_N,
 \qquad r_N=2K\sum_{n>N}\Lambda(n)n^{3/2}e^{-\pi n}\longrightarrow0.
\end{equation}
For every $u$ of finite weighted norm the series~\eqref{eq:physicalprime} converges absolutely and
\begin{align}
 p[u]&=\ip{K_P(u/\beta)}{u/\beta}_2,\label{eq:prepresentation}\\
 |p[u]-p_N[u]|&\le r_N\norm{u}_\omega^2,\qquad
 |p_N[u]|\le2\sum_{2\le n\le N}a_n\norm{u}_2^2.\label{eq:headtail}
\end{align}
\end{lemma}
\begin{proof}
Set $z=|x|+|x+s|\ge s$. Equation~\eqref{eq:envelopeW} and the arithmetic--geometric mean inequality give
\begin{align*}
 \beta(x)\beta(x+s)
 &\le K\exp\!\left(2z-\frac\pi2(e^{2|x|}+e^{2|x+s|})\right)\\
 &\le K e^{2z-\pi e^z}\le K e^{2s-\pi e^s}.
\end{align*}
The last step uses $2-\pi e^z<0$ for $z\ge0$. The norm of each weighted translation is bounded by the supremum of its multiplier, so its two-shift contribution at $s_n$ has norm at most $2Ka_nn^2e^{-\pi n}$. Since $\Lambda(n)\le\log n$, the norm series converges. Adjointness interchanges the two shifts, and reflection also interchanges them. This proves self-adjointness and even invariance.

Writing $u=\beta F$, the $n$th quadratic term equals the corresponding term in~\eqref{eq:physicalprime}. Its absolute value is bounded by the summable operator-norm bound times $\norm{F}_2^2=\norm{u}_\omega^2$. This proves both the identification and absolute convergence. The tail estimate follows from~\eqref{eq:rN}; Cauchy--Schwarz and translation invariance of ordinary $L^2$ give the finite-head estimate.
\end{proof}

\begin{remark}
Lemma~\ref{lem:prime} does not assert that $K_P$ is compact. Multiplication by a decaying function followed by a translation need not be compact on a nonatomic $L^2$ space. Nor does~\eqref{eq:rN} assert a uniformly vanishing tail for the unweighted incoming prime rate as the physical point escapes to infinity. The norm, the weight, and the coordinates in the lemma are essential.
\end{remark}

\section{The exact numerator domain}\label{sec:domain}
Define
\begin{equation}\label{eq:X}
 \begin{gathered}
 X=\{u\in\Le:\tg[u]<\infty,\ \norm u_\omega<\infty\},\\
 \norm u_\omega^2=\int_\R |u|^2/W\,dx,\qquad
 \norm u_X^2=\tg[u]+\norm u_\omega^2.
 \end{gathered}
\end{equation}

We first analyze $X$ directly from its two defining nonnegative terms. Because $W$ is positive and bounded,
\begin{equation}\label{eq:L2bound}
 \norm{u}_2^2\le\norm{W}_\infty\norm{u}_\omega^2.
\end{equation}
The Gamma form is closed on ordinary $L^2$, for example by the same difference-and-Fatou argument used for $\E$.

\begin{lemma}\label{lem:Xcore}
The space $X$ of~\eqref{eq:X} is Hilbert, and even $C_c^\infty(\R)$ is dense in $X$.
\end{lemma}
\begin{proof}
A sequence Cauchy in $X$ has a weighted $L^2$ limit and an ordinary Gamma-form limit. Equation~\eqref{eq:L2bound} identifies their ordinary $L^2$ limits, proving completeness.

Use the cutoffs $\chi_R$ from Proposition~\ref{prop:core}. For $h_R=(1-\chi_R)u$ the product identity gives
\begin{align*}
 |h_R(x+s)-h_R(x)|^2
 &\le2(1-\chi_R(x+s))^2|u(x+s)-u(x)|^2\\
 &\quad+2|u(x)|^2|\chi_R(x+s)-\chi_R(x)|^2.
\end{align*}
The first term tends to zero in the Gamma integral by dominated convergence. The second is dominated by
$2|u(x)|^2\min\{1,L^2s^2\}$, which is integrable by~\eqref{eq:rhointegrable}. Weighted convergence is another application of dominated convergence. Thus compactly supported vectors are dense in $X$.

For a fixed compactly supported vector, even smooth mollification converges in ordinary $L^2$ and in $\tg$, by the translation argument in Proposition~\ref{prop:core}. The supports remain in one compact interval, where $1/W$ is bounded. Ordinary $L^2$ convergence consequently gives weighted convergence as well.
\end{proof}

For original profiles and physical numerators the coordinates are
\begin{equation}\label{eq:coordinates}
 u=\kappa f,\qquad v=2cf=\frac uW,\qquad
 F=\frac u{\sqrt W}=\sqrt{2c\kappa}\,f.
\end{equation}
Only $f\mapsto F$ is the displayed unitary map from $\He$ to ordinary $\Le$. The map $f\mapsto u$ is an isometry to the weighted space with norm $\norm{\cdot}_\omega$, since
\begin{equation}\label{eq:weightedidentity}
 \norm{\kappa f}_\omega=\norm{f}_\mu.
\end{equation}

\begin{proposition}[Extension of the compact bridge]\label{prop:domain}
Multiplication by $\kappa$ gives a bijection $\D\to X$, with equivalent graph norms, and
\begin{equation}\label{eq:extendedbridge}
 q_0(\kappa f,\kappa g)=\E(f,g)-\tfrac12\ip{f}{g}_\mu,
 \qquad f,g\in\D.
\end{equation}
\end{proposition}
\begin{proof}
Put $C_0=|\cg|\norm{W}_\infty+\norm{K_P}$. On the compact smooth core, Theorem~\ref{thm:compactbridge}, Lemma~\ref{lem:prime}, and~\eqref{eq:weightedidentity} imply
\begin{align}
 \tg[\kappa f]&\le\E[f]+(C_0+\tfrac12)\norm{f}_\mu^2,\label{eq:graph1}\\
 \E[f]&\le\tg[\kappa f]+(C_0+\tfrac12)\norm{f}_\mu^2.\label{eq:graph2}
\end{align}
Taking compact-profile approximants from Proposition~\ref{prop:core} and using completeness of $X$ proves that $\kappa\D\subset X$.

Conversely, for $u\in X$ choose even $u_j\in C_c^\infty$ converging in $X$, by Lemma~\ref{lem:Xcore}. Since $\kappa$ is smooth and strictly positive, $f_j=u_j/\kappa$ are legitimate even compact smooth profiles. Equation~\eqref{eq:graph2} makes $(f_j)$ Cauchy in the original energy graph norm. Its weighted limit is $f=u/\kappa$, so closedness of $\E$ gives $f\in\D$. This proves surjectivity and graph-norm equivalence. The Gamma form and the prime form are continuous in $X$, so the quadratic identity passes to the limit; polarization gives~\eqref{eq:extendedbridge}.
\end{proof}

Recall $m(u)=\int_\R c(x)u(x)\,dx$. The integral exists absolutely for finite weighted norm, because
\begin{equation}\label{eq:meanbound}
 |m(u)|\le\left(\int_\R c^2W\,dx\right)^{1/2}\norm{u}_\omega
       =\tfrac12\norm{u}_\omega.
\end{equation}
Here $\int c^2W=\tfrac12\int c\kappa=1/4$. Thus $\mu(f)=2m(\kappa f)$, and the full pole form obeys
\begin{equation}\label{eq:polebridge}
 q[u]:=q_0[u]+2|m(u)|^2
       =\E[f]-\tfrac12\Var_\mu(f),\qquad u=\kappa f\in X.
\end{equation}
The coefficient $2$ in the pole is fixed by normalization. It is not an additional positive contribution available after centering.


\section{Sharp cuts and the centered exterior minimum domain}\label{app:cuts}
We record the domain argument, rather than identifying a minimum
closure with an arbitrary maximal finite-energy class.
Write $W_\R(h)=\int_\R |h|^2/\alpha^2$. Under $h=\kappa f$ it equals
$\norm f_{L^2(\mu)}^2$. Proposition~\ref{prop:domain}
shows that the original topology and
$(t_\Gamma(h)+W_\R(h))^{1/2}$ are equivalent on numerators.

For a smooth numerator $h$, expansion gives
\[
 t_\Gamma(h1_I)+t_\Gamma(h1_O)
 =t_\Gamma(h)+2\ip{h1_I}{\J_\Gamma(h1_O)}.
\]
For $0<s\le2$, \eqref{eq:rhoelementary} gives
$\rho(s)\le1/(2s)+1/4\le1/s$. For $s\ge2$, the function
$se^{-s/2}$ decreases, so
$s\rho(s)\le(2/e)/(1-e^{-4})<1$; the last inequality follows already
from $e>8/3$. Thus the unweighted Carleman bound for the Gamma crossing
is finite: $\rho(s)\le1/s$ and the two boundary Carleman maps give,
for example, $\norm{\J_\Gamma}_{L^2(O)\to L^2(I)}\le2\pi$.
Thus the right-hand correction is bounded by a constant times
$\norm h_2^2$, and $W_\R$ splits exactly. Since $\alpha^2$ is bounded,
both cut maps extend continuously in the numerator form topology.

For a fixed smooth compact numerator, approximate either step cut from
its own side by smooth killed-core cutoffs. The error $e_\delta$ is
supported in strips of width $\delta$ at $\pm B$ and satisfies
\[
 \norm{e_\delta}_\infty+\operatorname{TV}(e_\delta)\le C,
 \qquad\norm{e_\delta}_2^2+W_\R(e_\delta)\le C\delta.
\]
The bounded-variation translation estimate gives
\[
 \norm{e_\delta(\cdot+s)-e_\delta}_2^2\le C\min(s,\delta).
\]
Integrating against $\rho(s)=O(1/s)$ near zero and its integrable
large-shift tail proves
$t_\Gamma(e_\delta)=O(\delta(1+|\log\delta|))\to0$.
Consequently the cuts are in the corresponding killed minimum closures.
Apply continuity to a general minimum-domain numerator to obtain the
same conclusion. On the bounded local interval $W_\R$ is equivalent
to physical $L^2$, because $\alpha$ is continuous and strictly positive.

The constant function belongs to $\D$. Indeed even smooth cutoffs
$\chi_R$ of the constant satisfy $\kappa\chi_R\to\kappa$ in physical
$H^1$ and in $W_\R$ by theta decay. The topology equivalence and
$t_\Gamma(h)\le C\norm h_{H^1}^2$ make them Cauchy in the original
energy topology. Subtracting a constant is an admissible gauge change,
preserving energy and variance.

For an exterior limit with zero moment, choose even compact smooth
$\psi$ in $O$ with $m_O(\psi)=1$. If $u_n$ are killed-core approximants,
replace them by $u_n-m_O(u_n)\psi$. Moment continuity makes the
correction vanish in form norm. Hence the centered cut lies in
precisely $\D_O^0$, not just in the moment kernel of a larger domain.


\section{The exact centered-row quantities}\label{app:rows}
This appendix fixes the analytic interpretation of contract (b).
Use $K=64$, $J_0=512$, $d_0=6$ and $X=2$.
Let $\varphi_0=B^{-1/2}$ and
$\varphi_m=(2/B)^{1/2}\cos(k_mx)$ for $m>0$, where $k_m=m\pi/B$.
They are the unitary folds of the physical cosines.
All rows below are formed from the same matrix $Z$ and the same fixed rational vector $a$ in Hypothesis~\ref{hyp:response}.
The displayed formulas define quantities to be enclosed; their small numerical bounds remain assumptions.
For the physical columns $w_i$ of $\V$, let $Q_i=\J^*w_i$ denote
their weighted-dual kernels. This notation means
$\ip{\J u}{w_i}=\int_OuQ_i$ and
$\alpha Q_i\in L^2(O)$; it does not assert $Q_i\in L^2(O,\dd x)$.
If $m_O(u)=0$, then for every fixed vector $a$,
\begin{equation}\label{eq:centering}
 \V^*\J u=\int_Ou(x)[Q(x)-ac(x)]\dd x,\qquad
 \norm{\V^*\J u}\le\sqrt{\wo(u)}\norm{\alpha(Q-ac)}_{L^2(O;\R^{64})}.
\end{equation}
This is exact moment cancellation, not an approximate mean repair.
The choice of $a$ need not be accurate as a numerical minimizer.
The direction $c$ is subtracted in the weighted dual, not by physical
$L^2$ projection; $c$ is not a physical exterior $L^2$ vector.

Split the trial response into cosine rows below $K=64$ and rows
64 through 1023. Duality and Theorem~\ref{thm:source} give
\begin{equation}\label{eq:highrow}
 \norm{\alpha\J^*\V_{\ge64}}_{L^2(O;\R^{64})}
 \le j\norm{Z_{\ge64}}_{\rm HS}.
\end{equation}
This pays every discarded row of the finite trial response. The error
of the trial response as an approximation to $H_{64}^{-1}E$ belongs instead
to \eqref{eq:contracta}; it is not silently removed.

Let $b_m=\varphi_m(B)$ and $\lambda_j=2j+1/2$.
For $x>B$, direct cosine integration of the two folded exponential
kernels gives the low-response Gamma adjoint
\begin{equation}\label{eq:Qgamma}
 Q_{\Gamma,<K}(x)=\sum_{j\ge0}g_je^{-\lambda_j(x-B)},\qquad
 g_j=\frac{1-e^{-2\lambda_jB}}{\sqrt2}
       \sum_{m<K}\frac{b_m\lambda_j}{\lambda_j^2+k_m^2}Z_m.
\end{equation}
The factor $1/\sqrt2$ converts folded coefficients to physical
full-line columns. Put $A_J=2J_0+1/2$, and define signed moments
\[
 M_r=\sum_{m<K}b_m(-k_m^2)^rZ_m\quad(0\le r<d_0),\qquad
 H_{d_0}=\sum_{m<K}|b_m|\norm{Z_m}k_m^{2d_0}.
\]
The exact finite geometric identity for
$\lambda/(\lambda^2+k^2)$ has remainder of modulus at most
$k^{2d_0}/\lambda^{2d_0+1}$, without requiring $k<\lambda$.
Therefore
\[
 \norm{\sum_{m<K}\frac{b_m\lambda}{\lambda^2+k_m^2}Z_m}
 \le\sum_{r<d_0}\frac{\norm{M_r}}{\lambda^{2r+1}}
       +\frac{H_{d_0}}{\lambda^{2d_0+1}}.
\]
With $\sigma_p(A)=A^{-p}+A^{1-p}/[2(p-1)]$, decreasing integral
comparison gives $\sum_{j\ge J_0}\lambda_j^{-p}\le\sigma_p(A_J)$.
Minkowski's inequality and
$\norm{e^{-\lambda h}}_{L^2(0,\infty)}=(2\lambda)^{-1/2}$ prove
\begin{equation}\label{eq:egamma}
 e_\Gamma=\sqrt{A/2}\left[
 \sum_{r<d_0}\norm{M_r}\sigma_{2r+3/2}(A_J)
 +H_{d_0}\sigma_{2d_0+3/2}(A_J)\right]
\end{equation}
as a bound for the complete omitted weighted Gamma norm. Both
exterior halves are included. The signed moments are summed before
their vector norms are taken.

The remaining finite kernel is
\begin{equation}\label{eq:K0}
 \begin{split}
 K_0(x)={}&\sum_{j<J_0}g_je^{-\lambda_j(x-B)}-ac(x)\\
 &+\sum_{\substack{2\le n\le16\\|x-\log n|<B}}
  \frac{a_n}{\sqrt2}\sum_{m<K}|b_m|Z_m\cos(k_m(x-\log n)).
 \end{split}
\end{equation}
All ten prime powers through 16 are included:
$2,3,4,5,7,8,9,11,13,16$. Their Gamma, prime and mean contributions
are combined before vector norms. Boundary values at support switches
have measure zero.

The rough bound
\begin{equation}\label{eq:roughalpha}
 \alpha(x)^2\le80e^{4|x|-\pi e^{2|x|}}\quad(x\in O)
\end{equation}
follows already from the envelope: $4\pi^2<40$ and the tail-to-leading
coarse ratio $32e^{-3t}<1$ permit the constant 80.
Define
\[
 C_X=\sum_{j<J_0}\norm{g_j}
 +\left(\sum_{n\le16}a_n\right)
      \sum_{m<K}\frac{|b_m|\norm{Z_m}}{\sqrt2}+\norm a.
\]
For $x\ge X=2$, $\norm{K_0(x)}\le C_Xe^{x/2}$, including any finite
prime supports continuing past $X$. Since
$5-2\pi e^{2x}\le5-2\pi e^{2X}<0$, the full spatial tail is bounded by
\begin{equation}\label{eq:eX}
 e_X^2=\frac{160C_X^2e^{5X-\pi e^{2X}}}{2\pi e^{2X}-5}.
\end{equation}
The factor 160 includes both exterior half-lines.

On $[B,2]$ the nine support switches, in increasing order, are
\[
 \log5-B,\ \log7-B,\ \log8-B,\ \log9-B,\ \log2+B,
 \ \log11-B,\ \log13-B,\ \log3+B,\ \log16-B.
\]
The ten intervals have the fixed subdivision counts
\[
 (5,169,67,59,48,53,84,67,37,14),
\]
totaling 603.
On an exact coarse interval $[L,R]$ with count $N$, the exact finer
endpoints are $l_i=L+(R-L)i/N$, $r_i=L+(R-L)(i+1)/N$.
This algebraic identity proves adjacency and coverage. Overlap of
arbitrary interval boxes alone would not do so. Each open cell lies
between consecutive support switches and has a constant active-prime
set.

For a cell $[l,r]$, center $x_0$ and half-width $h$, define
$v_s=K_0^{(s)}(x_0)/s!$, $0\le s<16$. Taylor's theorem gives the
following explicit uniform vector remainder bounds:
\begin{align*}
 R_\Gamma&=\sum_{j<J_0}\norm{g_j}e^{-\lambda_j(l-B)}
                   \frac{(\lambda_jh)^{16}}{16!},\\
 R_p&=\left(\sum_{n\ {\rm active}}a_n\right)
       \sum_{m<K}\frac{|b_m|\norm{Z_m}}{\sqrt2}
                   \frac{(k_mh)^{16}}{16!},\\
 R_c&=\norm a\cosh(r/2)\frac{(h/2)^{16}}{16!}.
\end{align*}
Thus $M_{[l,r]}=\sum_{s=0}^{15}\norm{v_s}h^s+R_\Gamma+R_p+R_c$
bounds $\norm{K_0(x)}$ on the whole cell, and
\begin{equation}\label{eq:Ifin}
 I_{\rm fin}=\sum_{[l,r]}2(r-l)A(l)M_{[l,r]}^2
\end{equation}
bounds its full even weighted integral up to $X$.
Contract (b) requires outward enclosures of these
quantities; the analytical formula is a supremum-and-remainder bound,
not a quadrature accuracy inferred from mesh size. The subdivision counts do not replace these exact endpoint definitions.

\subsection{The complete prime tail in the response budget}
For $y\in I$, $|y\pm\log n|\ge\log n-B$ and each of the two
translated restrictions has weighted norm at most
$\sqrt{80}e^{-2B}n^2e^{-\zeta n^2}$, with
$\zeta=(\pi/2)e^{-2B}>3/10$. For $N\ge16$, integration by parts and
$\int_N^\infty e^{-\zeta x^2}\dd x\le
 e^{-\zeta N^2}/(2\zeta N)$ give
\[
 \norm{\J_{p,>N}}_{\wo\to L^2(I)}
 \le2\sqrt{80}e^{-2B}
 \left(\frac{N}{2\zeta}+\frac1{4\zeta^2N}\right)e^{-\zeta N^2}.
\]
At $N=16$, $2\sqrt{80}e^{-2B}<18$, the parenthesis is less than 27,
and $\zeta N^2>75$; hence the bound is smaller than
$b_{16}=500/20^{25}$. These entire integer tails dominate all omitted
prime powers and retain both orientations.


The finite part has weighted norm at most $\sqrt{I_{\rm fin}}+e_X$,
the complete omitted Gamma part has norm at most $e_\Gamma$,
and the complete omitted prime part has norm at most
$b_{16}\norm{Z_{<64}}_{\HS}$.
The high trial rows have norm at most $j\norm{Z_{64:1023}}_{\HS}$ by
\eqref{eq:highrow}. Exact exterior centering therefore gives
\[
 \norm{V^*Ju}\le
 \left(\sqrt{I_{\rm fin}}+e_X+e_\Gamma+
 b_{16}\norm{Z_{<64}}_{\HS}
 +j\norm{Z_{64:1023}}_{\HS}\right)\sqrt{\wo(u)}.
\]
This proves the decomposition used in \eqref{eq:responsebound};
the numerical caps themselves remain contract (b).

\section{Computational input identity and verification boundary}\label{app:provenance}
The proof has two different computational input classes.
Hypothesis~\ref{hyp:local} is an accepted whole-domain local theorem at $B=4/5$.
Its supplied packet gives $q_I[u]\ge10^{-22}\norm u_2^2$.
The records below identify that input and its computational evidence.
They are data-provenance identifiers, not citations to an earlier manuscript as a substitute for an analytic proof.
A cryptographic identity does not establish an inequality.

\begingroup\small
\begin{longtable}{@{}p{0.96\textwidth}@{}}
\toprule Local input artifact and SHA-256\\\midrule\endhead
Local theorem source\\*
\nolinkurl{ac04a393e8b6112d81bfb598c2dd00450b2a22eca36c298e3bbe82f7ef78b4c6}\\[3pt]
Local certificate archive version 1.1\\*
\nolinkurl{b805a431c4c77d955d5a6d4ed7397a01faf12cd727c1051af607ed14f77ab74b}\\[3pt]
Radius-four-fifths source diagnostic\\*
\nolinkurl{4c67f0c3f37bb8ce3c4f7c6da166050f2a1760c62bd47227c7709fb3a0b2ea05}\\[3pt]
Radius-four-fifths producer success record\\*
\nolinkurl{19156f6635bb25e853b9a1aecd650672043113f85d279c562e1a090269b0ca91}\\[3pt]
Exact-rational replay source\\*
\nolinkurl{3c5e006df82fbc985ce9b7ed73e26d62a9679086cc399dcb022224325c9fde13}\\[3pt]
Exact-rational replay result\\*
\nolinkurl{f2d4eef6e3efcd632ee252b2cca56be120e5ce5b3c59cdcc499277bc4ee528e2}\\[3pt]
Separate MPFR input archive\\*
\nolinkurl{4e57e261db7c8cfebdc5dc82b23d8aecfab3b639aa7fa5c95901b90f3aa18f8b}\\[3pt]
Separate MPFR evidence\\*
\nolinkurl{2a6fe3e8dfa04308ed555271bb19441b7d3ce783b6cc427419840a2d39a231d8}\\
\bottomrule
\end{longtable}\endgroup
The full local minimum domain, ordinary physical $L^2$ norm, positive mean pole, exact local prime powers $2,3,4$, and complete Gamma remainder are part of the input statement.
Finite matrix positivity alone is not the input theorem.
The supplied exact-rational replay and separately supplied source-enclosure evidence represent different validation layers.
No source generation, quadrature, local certificate, rational replay, or numerical sign gate was executed in preparing the present analytic manuscript.

The second class consists of Hypothesis~\ref{hyp:response}(a),(b).
One fixed data source for those contracts is the archive
\texttt{centered\_exterior\_cost\_sources.zip}, SHA-256
\begin{center}\small
\nolinkurl{4984450bc9f3174d354fdfe9731919dfb4c134f2db1ef0b55c40bc701b04792e}.
\end{center}
The relevant member identifiers are
\begingroup\small
\begin{longtable}{@{}p{0.96\textwidth}@{}}
\toprule Response input member and SHA-256\\\midrule\endhead
\texttt{centering\_vector.json}\\*
\nolinkurl{c04a7d7ac467971ae101290adbaec719826af343d7c0f589bdb7789c8ca8b4ed}\\[3pt]
\texttt{run/evidence.json}\\*
\nolinkurl{6c404a912f5986ae08080f2811a5d797d9aed22fc5683e24c9f2241525e4d3cd}\\[3pt]
\texttt{run/cells.json}\\*
\nolinkurl{d93657d28a399250f87f0d12f5421ced5c72178ff874c67ec5a3087a387ea6c7}\\[3pt]
\texttt{local\_response/run1024/response\_coefficients.json}\\*
\nolinkurl{b40fff3249a60dbb1224c9099360c96333d19fe22b7bd10bda126e9ece44f7a9}\\[3pt]
\texttt{local\_response/run1024/exact\_matrices.json}\\*
\nolinkurl{792718e3d299d74ca5757e5f17942608ab005468759673bd2f36388cf102dfba}\\[3pt]
\texttt{local\_response/run1024/evidence.json}\\*
\nolinkurl{62232dc8890083bfd967ad7fa3e854fe4ca00925b7b493a45dc7a45ecf234cbb}\\
\bottomrule
\end{longtable}\endgroup
The archive also identifies 60 finite local Gram blocks covering rows 1024 through 16383 and the infinite-tail enclosure.
Its local parameters are 1024 trial modes, cutoff 16384, 192-bit precision, 64 Gamma labels and tail order six.
The exterior parameters are fixed in Appendix~\ref{app:rows}.
These identify the declared inputs, not a new certification.
File identity, finite rational replay, and full source-enclosure validation remain distinct.
The clean source distribution contains no certificate data or executable research code; its build compiles the document only.


\section*{License and status}
Copyright \textcopyright\ 2026 Jongmin Choi.
This manuscript is licensed under Creative Commons Attribution--NonCommercial 4.0 International
(CC BY-NC 4.0): \url{https://creativecommons.org/licenses/by-nc/4.0/}.
The license does not imply external validation of the mathematical claims.
This version is a research manuscript; it has not been represented as journal accepted, peer reviewed, or formally proof verified.

\begingroup\small
\begin{thebibliography}{9}
\bibitem{DLMF}
F.~W.~J. Olver et al., eds., \emph{NIST Digital Library of Mathematical Functions}.
Jacobi transformation, Eq.~20.7.32, and digamma integral, Eq.~5.9.16.
\url{https://dlmf.nist.gov/20.7.E32};
\url{https://dlmf.nist.gov/5.9.E16}.
\bibitem{Garrett}
P. Garrett, \emph{Guinand's explicit formula},
University of Minnesota lecture notes, February 12, 2021, Theorem~0.1.
\newline
\href{https://www-users.cse.umn.edu/~garrett/m/complex/notes_2020-21/12g_guinand_explicit_fml.pdf}{University of Minnesota source PDF}.
\bibitem{GarrettZeros}
P. Garrett, \emph{Counting zeros of $\zeta(s)$},
University of Minnesota lecture notes, March 12, 2023, Theorem~0.1.
\newline
\href{https://www-users.cse.umn.edu/~garrett/m/complex/notes_2021-22/12c_counting_zeros.pdf}{University of Minnesota source PDF}.
\bibitem{Kato}
T. Kato, \emph{Perturbation Theory for Linear Operators},
second edition, Classics in Mathematics, Springer, Berlin, 1995, Chapter~VI.
\url{https://doi.org/10.1007/978-3-642-66282-9}.
\end{thebibliography}\endgroup
\end{document}
